\input{preamble}
\externaldocument[mumford-tate-]%
{mumford-tate-groups-in-hodge-theory}

\begin{document}

\title{Complex Geometry}
\maketitle

\phantomsection
\label{section-phantom}

\tableofcontents

\section{Complex structures on vector spaces}
\label{section-complex-structure}

\noindent
Let's discuss complex structures at the level
of linear algebra.

\begin{definition}
\label{definition-compl-struc-vecto-space}
An {\it almost complex structure} on a vector
space $V$
is an endomorphism $I \in \text{End}(V)$ such
that $I^2 = -\text{Id}_V$.
\end{definition}

\begin{lemma}
\label{lemma-almos-compl-struc-then-compl}
A vector space endowed with an almost complex
structure
admits a natural structure of a complex
vector space.
\end{lemma}

\begin{proof}
$(a+ib)v:=av + b I(v)$.
\end{proof}

\noindent
Conversely, given a complex vector space,
multiplication by $i$ produces an almost
complex structure.
Thus the two notions become the same.

\medskip\noindent
Now we introduce the complexification of a
real vector space.
Let $V_{\mathbb{C}}:=V
\otimes_{\mathbb{R}}\mathbb{C}$.
This means that elements of $V_{\mathbb{C}}$
are finite sums of elements of the form $v
\otimes (a+iv)$
for $a,b \in \mathbb{R}$ and $i = \sqrt{-1}$.

Notice that although $(V,I)$ is also a
complex vector space,
it is a different one from $V_\mathbb{C}$.
The complexification $V_\mathbb{C}$
is much more interesting since it gives us
the $(p,q)$-decomosition
when we take the exterior algebra.

Further, $V_\mathbb{C}=V \oplus i V$.
We define {\it conjugation} of vectors in
$V_{\mathbb{C}}$ by
$\bar{v}=\overline{v \otimes \lambda} := v
\otimes \bar{\lambda}$

This means that $\overline{V}$ is just $V$
but
with the scalar product $\lambda
v:=\bar{\lambda}v$.
for any $v \in V_{\mathbb{C}}:=V
\otimes_{\mathbb{R}}\mathbb{C}$.
This makes the equation
$\overline{\Lambda^{1,0}(V)}=\Lambda^{0,1}(V)
$,
which appears in the next proposition,
obvious.
(See \cite[p. 25]{huc}.)

\begin{lemma}
\label{lemma-i-eigenspace-decomposition}
We have a decomposition
$$
V_\mathbb{C}=V^{1,0}\oplus V^{0,1}
$$
given by the eigenspaces of $I$.
$\overline{\Lambda^{1,0}}=\Lambda^{0,1}$
holds by definition of conjugation.
\end{lemma}


Now let's discuss the $p,q$ decomposition.
The upshot is the following:
$$
T^2(V \oplus  W)=
T^2(V)\otimes T^0(W)
\oplus T^1(V)\otimes T^1(W)
\oplus T^0(V)\otimes T^2(W).
$$
Explicitly, the decomposition of a general
element on the LHS is
\begin{align*}
(v_1,w_1)\otimes(v_2,w_2)
&=(v_1,0)\otimes(v_2,w_2)+(0,w_1)\otimes(v_2,
w_2)\\
&=(v_1,0)\otimes(v_2,0)+(v_1,0)\otimes(0,w_2)
\\
&+(0,w_1)\otimes(v_2,0)+(0,w_1)\otimes(0,w_2)
\end{align*}
which is valid because the tensor product of
a vector space
$T^2X$ is just the pairs $x \otimes y$
subject
to relations that $(x_1+x_2)\otimes y=x_1
\otimes y+x_2\otimes y$
and addition for $y$ and scalar product.
So we just identify elements
\begin{align*}
(v_1,0)\otimes(v_2,0)
&\leftrightarrow
v_1 \otimes v_2 \in T^2(V) \otimes
T^0(W):=T^{2,0}(V \oplus W)\\
(v_1,0)\otimes(0,w_2)
&\leftrightarrow
v_1 \otimes w_2 \in T^1(V) \otimes T^1(W):=
T^{1,1}(V\oplus W)\\
(0,w_1)\otimes(v_2,0)
&\leftrightarrow
w_1 \otimes v_2 \in T^1(W) \otimes  T^1(V)
\cong T^1(V) \otimes  T^1(W):=T^{1,1}(V\oplus
W)\\
(0,w_1)\otimes(0,w_2)
&\leftrightarrow
w_1\otimes w_2 \in T^0(V) \otimes
T^2(W):=T^{0,2}(V \oplus  W).
\end{align*}

\noindent
We have proved that any element in $T^2(V
\oplus W)$
can be written as an element in the vector
space
$T^2(V)\otimes T^0(W)
+ T^1(V)\otimes T^1(W)
+ T^0(V)\otimes T^2(W)$.
Since the intersection of the factors is
disjoint,
the sum can be taken direct.

\medskip\noindent
Notice that we have two copies of
$T^{1,1}(V\oplus W)$.
By the binomial theorem and isomorphisms
of the kind $A \otimes B \cong B \otimes A$,
we readily see that
 $$
T^k(V \oplus W)
\cong
\bigoplus_{p+q=k}\binom{k}{p}T^p(V)\otimes
T^q(W).
$$
Of course the binomial coefficient matters
little
when the underlying field is $\mathbb{C}$.

The next step is to take the quotient on the
LHS
by the relation $v \otimes v=0$,
i.e., the ideal generated by $v \otimes v$.
When we compute $(v_1,w_1)\otimes(v_1,w_1)$
as above,
this relation becomes $v_1 \otimes v_1=0$
and $w_1\otimes w_2=0$ on the first and last
components.
If we want the middle component to vanish as
well
we must ask that the isomorphism
$T^1(W) \otimes  T^1(V)\cong T^1(V) \otimes
T^1(W)$
exchanges sign. Then we obtain that
$$
\Lambda^2(V \oplus W)
=\Lambda^{2,0}(V\oplus W)
\oplus\Lambda^{1,1}(V\oplus W)
\oplus\Lambda^{0,2}(V \oplus W)
$$
by defining $\Lambda^{p,q}(V \oplus
W):=\Lambda^p(V)\otimes \Lambda^q(W)$.

To make sure that this generalizes to all
$p,q$
write, e.g. $p+q=3$,
\begin{align*}
(v_1,w_1)\otimes(v_1,w_1)\otimes(v_1,w_1)
&=(v_1,0)\otimes(v_1,0)\otimes(v_1,0)\\
&+(v_1,0)\otimes(v_1,0)\otimes(0,w_1)\\
&+…
\end{align*}

\noindent
so that imposing the relation $v \otimes v$
on $T^3(V \oplus W)$
implies imposing it on every factor.
[Further arguments needed here… but looks
like just boring computations…]

\medskip\noindent
The main result of Hodge theory is that
when we take $T^{1,0}(M)\oplus T^{0,1}(M)$,
the eigenspaces of the complex structure,
and pass to complex singular cohomology of a
Kähler manifold we get
\begin{equation}
\label{equation-hodge-theorem}
H^{k}(M,\mathbb{C})
=\bigoplus_{p+q=k}H^{p,q}(M)
\end{equation}

\noindent
where $H^{p,q}(M)$ is the set of differential
forms
represented by a closed form in
$\Lambda^{p,q}(M)$.

\medskip\noindent
It also happens that
$\overline{\Lambda^{p,q}(V)}=\Lambda^{q,p}(V)
$,
which is due to the fact that
$\overline{\Lambda^{1,0}}=\Lambda^{0,1}$.


\begin{exercise}
\label{exercise-1,1-forms-are-i-invariant}
Show that $(1,1)$-forms are invariant under
the action of
the complex structure $I$.
And conversely, a 2-form is of type $(1,1)$
only if it is $I$-invariant.
\end{exercise}

\begin{proof}
Let $\alpha\wedge\beta \in \Lambda^{1,1}(M)$.
Recall that
$\Lambda^{1,1}(M)=\Lambda^{1,0}(M)
\otimes\Lambda^{0,1}(M)$,
i.e. $\alpha \in \Lambda^{1,0}(M)$ and
$\beta \in \Lambda^{0,1}(M)$.
For any vectors $v,w$ tangent to some point
$p \in M$,
\begin{align*}
I(\alpha\wedge\beta)(v,w)
&=(\alpha\wedge\beta)(Iv,Iw)\\
&=\alpha(Iv)\beta(Iw)-\beta(Iv)\alpha(Iw)\\
&=-\sqrt{-1}\sqrt{-1}(\alpha(v)\beta(w)-
\beta(v)\alpha(w))\\
&=\alpha(v)\beta(w)-\beta(v)\alpha(w)\\
&=(\alpha\wedge\beta)(v,w).
\end{align*}

\noindent
The converse follows easily by the direct sum
splitting of $\Lambda^2(M)$ and noting that
a $(2,0)$ or $(0,2)$-form $\gamma$
satisfies $I\cdot\gamma(v,w)=-\gamma(v,w)$
in an analogous computation.
\end{proof}

\noindent
In particular this implies that $\omega$,
the fundamental form of the Riemannian metric
$g$
and the complex structure $I$ is a
$(1,1)$-form.
So this shows that Kähler classes are in
$H^{1,1}(M)$.






\noindent
Now we shall continue discussing complex
linear algebra up
to Lefschetz hard theorem
\ref{theorem-lefschetz-decomposition}
and the Hodge-Riemann relations
\ref{theorem-hodge-riemann-relation}.


\begin{definition}
\label{definition-lefschetz-operator}
Let $(V,g,I)$ be an almost complex vector
space
equipped with an hermitian metric $g$.
The {\it Lefschetz operator} is
\begin{align*}
L: \Lambda^\bullet(V_{\mathbb{C}}^\vee)
&\longrightarrow
\Lambda^\bullet(V_{\mathbb{C}}^\vee) \\
\alpha &\longmapsto \omega\wedge\alpha.
\end{align*}
\end{definition}

\noindent
Recall that a metric on a vector space
induces a metric on the dual vector space
by defining the dual basis of an orthonormal
basis
to be orthonormal.
And for higher exterior powers all we do is
declare that wedge products of elements of
the dual orthonormal basis
(in ascending order)
form an orthonormal basis; more precisely
if $e_1,…,e_n$ is orthonormal for $V$ then
the elements of the form
$e_{i_1}\wedge…\wedge e_{i_k}$ with
$i_1<…<i_k$
are an orthonormal basis for
$\Lambda^k(V^\vee)$.

\begin{exercise}
\label{exercise-hodge-star-existence}
Let $V$ be an oriented Riemannian vector
space with
orientation form $\text{Vol}$.
There is a unique operator
$*:\Lambda^\bullet(V^\vee) \to
\Lambda^\bullet(V^\vee)$
called the {\it Hodge star} such that
$$
\alpha\wedge*\beta=g(\alpha,\beta)\text{Vol}
$$
for any $\alpha,\beta \in \Lambda^k(V^\vee)$.
\end{exercise}

\begin{proof}
First prove uniqueness,
then prove that mapping a basic form
$e^{i_1}\wedge… \wedge e^{i_k}$
to $e^{j_1}\wedge…\wedge e^{j_{n-k}}$
where $(i_1,…,i_k,j_1,…,j_{n-k})$ is a
permutation of
$(1,…,n)$ will work (up to sign maybe).
\end{proof}

\noindent
Intuition on the Hodge star operator is that
it maps a form to its ``complement'' with
respect to the volume form.
For instance, $*dx=dy \wedge dz$ in
$\mathbb{R}^3$.

\begin{lemma}
\label{lemma-hodge-star-isometry}
Hodge star is an isometry.
\end{lemma}

\begin{proof}
Consider two basic covectors $e^I$ and $e^J$.
Then $*e^I$ and $*e^J$ are determined by $I$
and $J$.
By definition of the metric on
$\Lambda^\bullet(M)$,
$\left<e^I,e^J\right>$ is zero if $I\neq J$
and 1 otherwise.
Similarly, $\left<*e^I,*e^J\right>$ is zero
if $I\neq J$,
while, if $I=J$, then it will clearly be 1.
\end{proof}

\begin{lemma}
\label{lemma-hodge-star-squared}
\begin{reference}
\cite[Lemma 5.5]{voi}
\end{reference}
$$
*^2=(-1)^{k(n-k)}\text{Id}.
$$
\end{lemma}

\begin{proof}
We prove that
$$
\left<\alpha-*^2\alpha,\beta\right>=0\qquad
\forall \beta \in \Lambda^k(M).
$$
\begin{align*}
\left<\alpha,\beta\right>\text{Vol}
&=\left<*\alpha,*\beta\right>\text{Vol}\\
&=*\beta\wedge *^2\alpha\\
&=(-1)^{\text{deg}(*\alpha)
\text{deg}(\beta)}*^2\alpha\wedge *\beta\\
&=\left<(-1)^{(n-k)k}*^2\alpha,
*\beta\right>\text{Vol}.
\end{align*}
\end{proof}

\noindent
In the setting of linear algebra alone,
there is no confusion on which inner product
we use to define adjoint operators.

\begin{definition}
\label{definition-dual-lefschetz}
The adjoint operator of $L$ with respect to
the metric
is called {\it dual Lefschetz operator}, i.e.
it is
the unique operator $\Lambda$ such that
$$
\left<L \alpha,\beta\right>=\left<\alpha,
\Lambda \beta\right>.
$$
\end{definition}

\begin{exercise}
\label{exercise-dual-lefsc-hodge-twist}
$\Lambda=*^{-1}\Lambda *$.
\end{exercise}

\begin{proof}
\begin{align*}
\left<L \alpha,\beta\right>
&=L \alpha\wedge*\beta\\
&=\omega\wedge\alpha\wedge*\beta\\
&=\alpha\wedge\omega\wedge*\beta\\
&=\alpha\wedge L * \beta\\
&=\alpha\wedge * *^{-1}L * \beta\\
&=\left<\alpha,*^{-1}L *
\beta\right>\text{Vol}.
\end{align*}
\end{proof}

\noindent
Next we define $H$, another operator on
$\Lambda^\bullet$
called the {\it counting operator},
which is just the identity multipliied by the
codegree of a form:
$$
H|_{\Lambda^kV}:=(k-n)\text{id}.
$$
\begin{proposition}
\label{proposition-sl2-relat-exter-algeb}
Let $(V,\left<\cdot,\cdot\right>)$ be an
euclidean vector space
endowed with an compatible almost complex
structure $I$.
Then
$$
[H,L]=2L,\qquad [H,\Lambda]=-2\Lambda,\qquad
[L,\Lambda]=H,
$$
that is, there is a representation of
$\mathfrak{sl}_2$
on $\Lambda^\bullet(V)$.
\end{proposition}

\begin{definition}
\label{definition-primitive-form}
Let $(V,\left<\cdot,\cdot\right>)$ be an
euclidean
vector space with a compatible almost complex
structure $I$.
{\it Primitive forms} are the elements of
$\Ker \Lambda$.
The set of primitive forms of degree $k$
is denoted $P^k \subset \Lambda^k(V^\vee)$.
\end{definition}

\begin{theorem}[Lefschetz decomposition]
\label{theorem-lefschetz-decomposition}
$$
\Lambda^k(V^\vee)
=
\bigoplus_i L^i(P^{k-2i})
$$
holds and is called the {\it Lefschetz
decomposition}.
Moreover, the decomposition is orthogonal
with respect to
$\left<\cdot,\cdot\right>$.
\end{theorem}

\noindent
The upshot about this theorem
is that it is proved using the fact
that $\Lambda^\bullet(V^\vee)$ is an
$\mathfrak{sl}_2$-representation.

\begin{definition}
\label{definition-hodge-riemann-pairing}
\begin{reference}
\cite[Definition 1.2.35]{huc}
\end{reference}
Let $(V,\left<\cdot,\cdot\right>)$
be an euclidean vector space with a
compatible almost complex structure.
The {\it Hodge-Riemann pairing} is

\begin{align*}
Q: \Lambda^k(V^\vee)\times \Lambda^k(^\vee)
&\longrightarrow
\Lambda^{2n}(V^\vee)\cong \mathbb{R} \\
Q(\alpha,\beta)&
=(-1)^{\frac{k(k-1)}{2}}
\alpha\wedge\beta\wedge\omega^{n-k}.
\end{align*}
\end{definition}

\begin{theorem}[Hodge-Riemann bilinear
relation]
\label{theorem-hodge-riemann-relation}
$Q$ is positive definite on $P^{p,q}$, more
precisely,
$$
\sqrt{-1}^{p-q}Q(\alpha,\bar{\alpha})
=(n-(p+q))!\left<\alpha,
\alpha\right>_{\mathbb{C}}>0,
$$
for $\alpha \in P^{p,q}\setminus\{0\}$ and
$p+q\leq n$.
\end{theorem}

\medskip\noindent
Now let's discuss the local theory on complex
manifolds,
which is rather just analysis on an open set
of the complex numbers.
We start defining $\partial,\bar\partial$
and aim at proving Dolbeault theorem,
which is proved exactly in the same way as de
Rham's theorem.

\medskip\noindent
Upshot. We recall by sheaf theory
that the existence of an acyclic resolution
of a sheaf
implies that the cohomology of the resolution
equals the cohomology of the sheaf (see next
paragraph).
de Rham's complex is an acylic resolution
of the constant sheaf
$\underline{\mathbb{R}}$;
acyclicity is due to Poincaré lemma - a local
property.
There some way of viewing the complex of
singular cochains
as an acyclic resolution of the constant
sheaf $\underline{\mathbb{R}}$.
Thus both de Rham cohomology and singular
cohomology
are isomorphic to the cohomology of the
constant sheaf $\underline{\mathbb{R}}$.

The derived functors assign a resolution to
every object
in the abelian category if it has enough
injectives.
An acylcic resolution is such that the
derived functors
of all objects in the resolution vanish.
The cohomology of a sheaf is defined
as the resolution given by its exact
functors.

Dolbeault theorem is exactly the same:
Dolbeault's complex is acylcic by
$\bar\partial$-Poincaré lemma,
making it a resolution of the sheaf
$\Omega^p$.

There is an interesting relationship between
the two results,
which we anticipated since $\bar\partial$
appears to act on a ``piece'' of the exterior
algebra,
fixing another piece.

\medskip\noindent
Let us start by defining $\partial$ and
$\bar\partial$.
Take a real-valued smooth function $f \in
C^\infty(V,\mathbb{R})$.
We know that the de Rham differential is
given
in local coordinates $x^i,y^j$ of our $2n$
real-dimensional manifold by
$$
df=\partial_{x^i}fdx^i + \partial_{y^j}fdy^j.
$$
When we complexify the exterior algebra we
obtain
$\Lambda_{\mathbb{C}}^\bullet(V)
= \Lambda^\bullet(V)
\otimes_\mathbb{R}\mathbb{C}
= \Lambda^\bullet(T^*V
\otimes_\mathbb{R}\mathbb{C})$.
Here we can multiply by complex numbers,
so we can define
$$
\partial_z:= \frac{1}{2}(\partial_x  - i
\partial_y ),\qquad
\partial_{\bar{z}}:=\frac{1}{2}(\partial_x +
i \partial_y)
$$
and obtain
$$
dz=dx+idy,\qquad  d\bar{z}=dx - i dy.
$$

We see that
\begin{align*}
\partial_zfdz+\partial_{\bar{z}} fd\bar{z}
&=\frac{1}{2}(\partial_xf-i\partial_y
f)(dx+idy)\\
&+\frac{1}{2}(\partial_x f+i \partial_y
f)(dx- i dy)\\
&=\partial_x f dx + \partial_y f dy\\
&=df.
\end{align*}

\noindent
\begin{definition}
\label{definition-dbar}
\begin{reference}
\cite[Definition 1.3.5]{huc}
\end{reference}
Let $X$ be an almost complex manifold.
(Or vector space! The following
objects may be defined over $V_{\mathbb{C}}$;
we put almost complex manifold since the
definition
for almost complex manifolds shall be
identical)

If  $d:\mathcal{A}^k_{X,\mathbb{C}}\to
\mathcal{A}^{k+1}_{X,\mathbb{C}}$
is the $\mathbb{C}$-linear extension of the
exterior differential, then one defines
$$
\partial:=\Pi^{p+1,q} \circ d
:\mathcal{A}^{p,q}_X \to
\mathcal{A}^{p+1,q}_X,
\qquad
\bar{\partial}:=\Pi^{p,q+1} \circ d
:\mathcal{A}^{p,q}_X \to A^{p,q+1}_X
$$
where $\Pi^{p,q}$ is the projection.
\end{definition}

\noindent
From this definition we see that
$$
\partial f = \partial_z fdz,\qquad
\bar\partial f = \partial_{\bar{z}}fd\bar{z}
$$
for $f$ real or complex-valued.

If we take $f$ complex-valued
and put $f=u+iv$,
the condition $\bar\partial f =0$
becomes the Cauchy-Riemann equations.
Thus a complex-valued function is holomorphic
if and only if
$\bar\partial f=0$, which means that its
complexified de Rham differential is a
$(1,0)$-form,
i.e. a multiple of $dz$ only.

This way we see that since real partial
derivatives commute,
\begin{align*}
\partial_z \partial_{\bar{z}}
&=\left(\frac{1}{2}(\partial_x - i
\partial_y)\right)
\left(\frac{1}{2}(\partial_x+i
\partial_y)\right)\\
&=\frac{1}{4}(\partial_x^2
+i \partial_x\partial_y-i\partial_y\partial_x
-i^2 \partial_y^2)\\
&=\frac{1}{4}(\partial_x^2+\partial_y^2)\\
&=\frac{1}{4}(\partial_x^2
-i \partial_x\partial_y+i\partial_y\partial_x
-i^2 \partial_y^2)\\
&=\left(\frac{1}{2}(\partial_x + i
\partial_y)\right)
\left(\frac{1}{2}(\partial_x-i
\partial_y)\right)\\
&=\partial_{\bar{z}}\partial_z,
\end{align*}

\noindent
so that
\begin{align*}
\partial \bar\partial f
&=\partial(\partial_{\bar{z}}f d\bar{z})\\
&=\partial_z \partial_{\bar{z}}f dz \wedge
d\bar{z}.
\end{align*}

\noindent
And similarly,

\begin{align*}
\bar\partial \partial f
&=\partial_{\bar{z}}(\partial_z f dz)\\
&=\partial_{\bar{z}}\partial_zf d\bar{z}
\wedge dz\\
&=-\partial_z\partial_{\bar{z}}fdz\wedge
d\bar{z}
\end{align*}

\noindent
and thus
$$
\partial\bar\partial=-\bar\partial \partial
$$
on functions.

For an arbitrary $(p,q)$-form
$$
\alpha=f_{I,J}dz^I \wedge d\bar{z}^J,\qquad
|I|=p,|J|=q,
$$
we shall also have
$\partial\bar\partial=\bar\partial\partial$.



\medskip\noindent
We have the following local result:
\begin{lemma}[$\bar\partial$-Poincaré lemma
on the open disc]
\label{lemma-poincare-dbar-local-lemma}
\begin{reference}
\cite[Corollary 1.3.9]{huc}
\end{reference}
If $\alpha \in \mathcal{A}^{p,q}(B)$
is $\bar\partial$ closed, then there exists
$\beta \in \mathcal{A}^{p,q-1}$ such that
$\alpha=\bar\partial \beta$.
\end{lemma}

\begin{proof}
Looks like this is a consequence of complex
analysis.
\end{proof}

\begin{proposition}
\label{proposition-leibniz-for-del}
\begin{align*}
\partial(\alpha\wedge\beta)&=
\partial\alpha\wedge\beta
+(-1)^{p+q}\alpha\wedge\partial\beta\\
\bar\partial(\alpha\wedge\beta)
&=\partial\alpha\wedge\beta+(-1)^{p+q}
\alpha\wedge\bar\partial\beta.
\end{align*}
Also,
$$
\partial^2=\bar\partial^2=0,\qquad \partial
\bar\partial=-\bar\partial \partial.
$$

\end{proposition}


\section{Complex manifolds}
\label{section-complex-manifolds}

\noindent
Finally we arrive at complex manifolds.

\begin{definition}
\label{definition-complex-manifold}
A {\it complex manifold} $M$ is a smooth
manifold admitting an open cover
$\{U_\alpha\}$ and coordinate maps
$\varphi_\alpha:U_\alpha\to\mathbb{C}^n$ such
that $\varphi_\alpha\circ\varphi_\beta^{-1}$
is holomorphic on
$\varphi_\beta(U_\alpha\cap
U_\beta)\subset\mathbb{C}^n$ for all
$\alpha,\beta$.
\end{definition}

\begin{example}
\label{example-projective-space}
The standard coordinate charts of projective
space
$\mathbb{P}^n$ are
$$
U_i:=\{(z_0:…:z_n): z_i \neq 0\}
$$
along with the maps
\begin{align*}
\varphi_i: U_i &\longrightarrow \mathbb{C}^n
\\
(z_0:…:z_n) &\longmapsto
\left(\frac{z_0}{z_i},…,\frac{z_{i-1}}{z_i},
\frac{z_{i+1}}{z_i},…,\frac{z_n}{z_i}\right).
\end{align*}

\noindent
Transition functions are given by … which are
clearly holomorphic.
\end{example}

\noindent
As a warning note that we can take the charts
to be
$U_i=\{z_i=1\}$,
but when considering $U_i \cap U_j$
we must be careful since this is {\it not}
$\{z_i=1\}\cap \{z_j=1\}$.
For example in $\mathbb{P}^1$, this would say
that the intersection
$U_0 \cap U_1$
is the single point $(1:1)$, but what about
$(1:2)$?

\begin{definition}
\label{definition-almost-complex-manifold}
An {\it almost complex manifold} $M$ is a
smooth manifold together with
an endomorphism of the real tangent bundle
$I:TM \to TM$
such that $I^2=-\text{Id}$.
\end{definition}

\begin{proposition}
\label{proposition-complex-structure}
A complex manifold has a natural almost
complex structure.
\end{proposition}

\begin{proof}
Define bundle maps $I_i$ on local
trivializations $U_i$
by rotating vectors $90$ degrees in
counter-clockwise
direction at every copy of $\mathbb{C}$ in
the tangent space $\mathbb{C}^n$.
The upshot is that since $X$ is a complex
manifold
we can interpret the $I_i$ as multiplication
by $\sqrt{-1}$
in the complex vector space structure of the
tangent space.

Then the changes of coordinates satisfy the
condition
needed to glue to a global endomorphism,
namely
$I_j=g_{ij}I_ig_{ij}^{-1}$, trivially because
the $g_{ij}$ are holomorphic and thus
$\mathbb{C}$-linear.

\noindent
{\bf Addendum.} It looks like we can just
directly
define the $I_i$ by multiplication by
$\sqrt{-1}$.
\end{proof}


\noindent
But what about the converse?
What's the condition for an almost complex
manifold to be complex?

\begin{definition}
\label{definition-integ-almos-compl-struc}
An almost complex structure on a smooth
manifold $X$
is called {\it integrable} if it is induced
from a
complex structure on $X$ as in Proposition
\ref{proposition-complex-structure}.
\end{definition}

\begin{theorem}[Newlander-Nirenberg]
\label{theorem-newlander-nirenberg}
An almost complex structure $I$
on a smooth manifold is integrable if and
only if
$[T^{0,1},T^{0,1}]\subset T^{0,1}$.
\end{theorem}

Other equivalent conditions…

The {\it tensor} is given by
$$
N_I[X,Y]=[X,Y]+I[IX,Y]+I[X,IY]-[IX,IY]
$$




\medskip\noindent
Now let's discuss Dolbeault cohomology.


\begin{proposition}
\label{proposition-dolbeault-proposition}
\begin{reference}
\cite[Proposition 2.6.11]{huc}
\end{reference}
The space $H^0(X,\Omega^p_X)$ of holomorphic
$p$-forms on a complex manifold is the
subspace
$\{\alpha \in
\mathcal{A}^{p,0}(X):\bar\partial=0\}$.
\end{proposition}

\begin{proof}
The problem is local.
We use that $f dz_{i_1}\wedge…\wedge
dz_{i+p}$ is holomorphic
if and only if $f$ is holomorphic, which is
equivalent to
$$
\bar\partial(fdz_{i_1}\wedge…\wedge dz_{i_p}
=\sum \frac{\partial f}{\partial
\overline{z}_i}d\overline{z}_i
\wedge dz_{i_1}\wedge…\wedge dz_{i_p}=0.
$$
\end{proof}

\noindent
For a more general statement we must define
the {\it $(p,q)$-Dolbeault cohomology} by
 $$
H^{p,q}(X)_{\bar\partial}
:=\frac{\Ker
(\bar\partial:\mathcal{A}^{p,q}(X) \to
\mathcal{A}^{p,q+1}(X)}
{\text{Im}(\bar\partial:\mathcal{A}^{p,q-1}
(X)
\to \mathcal{A}^{p,q}(X))}.
$$
Then

\begin{theorem}[Dolbeault]
\label{theorem-dolbeault}
\begin{reference}
\cite[Corollary 2.6.21]{huc}, \cite[Theorem
4.2]{voi}
\end{reference}
The Dolbeault cohomology of $X$ computes the
cohomology
of the sheaf $\Omega^p_X$, i.e.
$H^{p,q}_{\bar\partial}(X)=H^q(X,\Omega^p_X)
$.
\end{theorem}



\section{Kähler manifolds, Kähler identities}
\label{section-kahler}


\begin{definition}
\label{definition-hermitian-structure}
A {\it hermitian structure} on a complex
manifold $M$
is a Riemannian metric $g$ that is also an
isometry for
the complex structure $I$, that is
$g(I\cdot,I\cdot)=g(\cdot,\cdot)$.
\end{definition}

\begin{definition}
\label{definition-fundamental-form}
If $(M,I,g)$ is a hermitian manifold,
$\omega(\cdot,\cdot):=g(I\cdot,\cdot)$
is called the {\it fundamental form}.
$(M,I,g)$ is called {\it Kähler} if its
fundamental form
is closed.
\end{definition}

\noindent
As in with vector spaces endowed with almost
complex structures,
we have the operators $L,\Lambda,*$ defined
on the
cotangent spaces of a complex manifold.
As a corollary we have that Lefschetz
decomposition
passes to exterior algebra
of a hermitian manifold:

\begin{lemma}
\label{lemma-lefsc-decom-manif}
Let $(X,g)$ be an hermitian manifold.
Then there exists a direct sum decomposition
of vector bundles
$$
\Lambda^k(X)=\bigoplus_{i\geq
0}L^i(P^{k-2i}X),
$$
there $P^{k-2i}X$ is the kernel of $\Lambda$.
\end{lemma}

\medskip\noindent
Apart from the linear-algebraic operators,
on manifolds we have other operators.

\begin{definition}
\label{definition-adjoi-de-rham-diffe}
The {\it adjoint de Rham differential} is
$$
d^*:=(-1)^{n(k+1)+1}*\circ d\circ *
:\mathcal{A}^k(M) \to \mathcal{A}^{k-1}(M)
$$
where, as always, $\mathcal{A}^k(M)$ denotes
the
sections of $\Lambda^k(M):=\Lambda^k(T^*M)$.
\end{definition}

\begin{definition}
\label{definition-laplace}
The {\it Laplace operator} is
$$
\Delta=d d^*+d^*d.
$$
\end{definition}

\noindent
If $M$ is a complex manifold,
then
$$
d^*=-*\circ d \circ *.
$$

Similarly we define
$$
\partial^*:=-* \circ \partial \circ *,
\qquad
\bar\partial^*:=-* \circ \bar\partial \circ
*,
$$
which yield the Laplacians
\begin{align*}
\Delta_\partial&=\partial \partial^*
+\partial^*\partial,\\
\Delta_{\bar\partial}&=\bar\partial
\bar\partial^*+\bar\partial^*\bar\partial.
\end{align*}

\begin{proposition}[Kähler identities]
\label{proposition-kahler-identities}
Let $X$ be a complex manifold endowed with a
Kähler metric $g$. Then
\begin{enumerate}
\item $[\bar\partial,L]=[\partial,L]=0$,
$[\bar\partial^*,\Lambda]=[\partial^*,
\Lambda]=0$.
\item  Etc.
\item
$\Delta_\partial=\Delta_{\bar\partial}=
\frac{1}{2}\Delta$.
\item $\Delta$ commutes with
$*,\partial,\bar\partial,\partial^*,
\bar\partial^*,L$
and $\Lambda$.
\end{enumerate}
\end{proposition}

\begin{definition}
\label{definition-dc}
On a complex manifold $(M,I)$,
$$
d^c:=I^{-1}d I=-I d I.
$$
\end{definition}

Notice (prove) that
$$
d^c=\frac{i}{2}(\bar\partial-\partial)
$$
so that
$$
dd^c=i(\partial \bar\partial).
$$
We may also define
$$
{d^c}^*:=-* \circ d^c \circ *.
$$




\noindent
Now let's introduce Fubini-Study form.
The following discussion of psh functions is
because
to pass from a potential, i.e. a function,
to a $(1,1)$-form we apply $dd^c$ (or
$\partial\bar\partial$).
Now if the potential is psh,
then the Hessian will be positive definite,
which is just the same as saying that
$dd^cf$ is positive definite $(1,1)$-form.

\begin{definition}
\label{definition-positive-11form}
Let $V$ be a real vector space equipped with
a complex structure $I$.
A $(1,1)$-form $\omega$ is {\it positive}
if $\omega(x,Ix)\geq 0$ for any nonzero $x
\in V$.
(Notice the $I$ is in the second entry!)
It is called {\it Hermitian} or {\it
semipositive}
if $\omega(x,Ix)>0$ for all nonzero $x \in
V$.
\end{definition}

\noindent
See \cite[Lemma 3.1.7]{huc}
to find Huybrecht's definition of a closed
{\it positive definite}
$(1,1)$-form, namely if
$\omega=\frac{i}{2}\sum h_{ij}dz_i \wedge
d\bar{z}_j$
such that $(h_{ij}(x))$ is a positive
definite hermitian matrix for any $x \in X$.

The point is that this condition
is the same as saying that the associated
Riemannian metric is in fact positive
definite.
Indeed, $\omega(\cdot,I\cdot)=g(\cdot,\cdot)$
which is positive definite.

Finally observe that numerical positivity of
the
Kähler class,
i.e. positivity with respect to the
intersection form,
follows exactly from this fact: we shall have
that
$\int_Y \omega^p>0$ for any analytic cycle
$Y$.

\begin{definition}
\label{definition-psh}
A function $f$ on a complex manifold is
{\it plurisubharmonic} if $dd^cf$ is
a positive $(1,1)$-form.
\end{definition}

\noindent
Again: psh is equivalent to the Hermitian
matrix
$$
\left(\frac{\partial^2 f}{\partial
z_i\partial \bar{z}_j}\right)
$$
being positive semidefinite.
ChatGPT: it is also equivalent to: its
restriction to
any complex line being subharmonic ($\Delta
f\geq 0$).

\begin{exercise}
\label{exercise-logl}
Let $z_1,…,z_n$ be the standard coordinates
on $\mathbb{C}^n$,
and $l(z_1,…,z_n):=\sum|z_i|^2$.
Prove that
$dd^c\text{log}l=\frac{dd^cl}{l}-\frac{dl
\wedge d^cl}{l^2}$
on $\Omega=\mathbb{C}^n\setminus 0$.
Prove that the Hermitian form $dd^c
\text{log} l$
is degenerate and can be written as
$(-\sqrt{-1})\sum_{i=1}^{n-1}
\alpha_i\xi_i\wedge \bar{\xi}_i$,
where $\xi_i$ is a local frame in
$\Lambda^{1,0}(\Omega)$
and $\alpha_i$ are positive functions.
\end{exercise}

\begin{proof}
For the first equation we just use the fact
(why?) that both $d^c$ and $d$ satisfy the
chain rule.
For the second equation I computed that
\begin{equation}
\label{equation-ddcl1}
dd^cl=\sqrt{-1}\sum dz_i\wedge d\bar{z}_i.
\end{equation}
But I also need to compute
$d^cl=\frac{\sqrt{-1}}{2}(\bar\partial-
\partial)l$
and $dl=(\partial+\bar\partial)l$.
I find that $dl \wedge d^cl=\sqrt{-1}\partial
l \wedge \bar\partial l$.
But I also found that
$\partial l=\sum \bar{z}_idz_i$ and
$\bar\partial l=\sum z_i d\bar{z}_i$.
So I get
\begin{equation}
\label{equation-ddcl2}
dl\wedge d^cl=\sqrt{-1}
\left(\sum \bar{z}_i dz_i\right)
\wedge
\left(\sum z_i d \bar{z}_i\right)
=\sqrt{-1}\sum\bar{z}_iz_jdz_i\wedge
d\bar{z}_j
\end{equation}

\noindent
Now we join Equations \ref{equation-ddcl1}
and \ref{equation-ddcl2}
to find the coefficient functions with
respect to $dz_i \wedge d\bar{z}_i$.
From Equation \ref{equation-ddcl1} we get
coefficient functions $\delta_{ij}$,
and from Equation \ref{equation-ddcl2} we get
$\bar{z}_iz_j$ so
$$
dd^c\text{log}l=\sqrt{-1}\sum
\left(\frac{\delta_{ij}}{l}
-\frac{\bar{z}_iz_j}{l^2}\right)
dz_i \wedge d\bar{z}_j.
$$
I'm not sure how to prove that  the
coefficient functions
are positive… Or that the form is degenerate…
But fortunately this can be checked at
\cite[p. 118]{huc}.
\end{proof}

\noindent
The point of this exercise is to show that
the form
$dd^c\ell$ is almost an Hermitian form ---
but it has a kernel.
I think the point of the next exercises in
this
\href{http://verbit.ru/IMPA/Surfaces-2025/%
assign-surfaces-2025-04.pdf}{
handout}
is to show that the kernel is actually
the ``radial directions''.
And the point of that is to take quotient
by $\mathbb{C}^*$ action, which is literally
passing to the projective space.

A bit more formally:

\begin{definition}
\label{definition-basic-form}
Let $M$ be a manifold and $\omega \in
\Omega^k(M)$.
Suppose that $\pi:M \to B$ is a surjective
submersion with connected fibers.
We say that $\omega$ is {\it basic} (with
respect
to $\pi$) is there exists a form
$\overline{\omega}\in \Omega^k(B)$ such that
$\pi^*  \overline{\omega}$.
\end{definition}

\begin{exercise}
\label{exercise-basic-forms-characterization}
\begin{enumerate}
\item Show that $\omega$ is basic if and only
if
$i_X \omega=0$ and $L_X\omega=0$
for all vector fields $X$ tangent
to the fibers of $\pi$.
In particular, if $\omega$ is closed,
show that it is basic if
$\Ker (\pi_*) \subseteq \Ker(\omega)$
(pointwise in $M$).

\item Suppose that $\omega$ is a closed
2-form on $M$ and $\Ker(\pi_*)=\Ker(\omega)$.
Show that $\omega=\pi^*\overline{\omega}$
and $\overline{\omega}\in\Omega^2(B)$
is symplectic.

\item (Application to reduction.)
\end{enumerate}
\end{exercise}

\begin{proof}
\begin{enumerate}
\item Since $\pi$ is surjective, for any
vector $\bar{X}$ tangent to $B$ there is
a vector $X$ tangent to $M$ such that
$\pi_*X=\bar{X}$.
Define a form
$\overline{\omega}(\bar{X}^1,…,\bar{X}^k):=
\omega(X^1,…,X^k)$.

To show this does not depend on the choice of
lift of $\bar{X}$
just notice that the tangent space of $M$ can
be written
$T_pM \simeq T_pF_b \oplus \pi^*(T_bB)$ where
$\pi(p)=b$.
Since $\omega$ kills vectors tangent to the
fibers,
$\overline{\omega}$ is well-defined.

Notice that there can be different choices of
$p$ such that
$\pi(p)=b$. We must show that the definition
is also
independent of this choice.
Let $X \in T_pM$ and $Y \in T_qM$ be such
that
$\pi(p)=\pi(q)=b \in B$ and $\pi_*X=\pi_*Y$.
We wish to show that
$\omega(…,X,…)=\omega(…,Y,…)$.

This a little more involved so I'll just
sketch it for now.
If there is a vector field $V$ tangent to the
fibers
whose flow $\varphi$ joins $p$ and $q$, we
are done,
since we pull back $\omega$ along this flow,
and by the condition of the Lie derivative,
i.e. $L_V\omega$ for $V$
tangent to the fiber, we see that
$\varphi^*\omega=\omega$.
(Indeed,
$L_V\omega=\frac{d}{dt}\varphi^*\omega$,
so the derivative is zero and at $t=0$ we
have $\omega$.)

Note that the choice of $Y$ within $T_qM$
is superfluous since we already know that
$\omega$
is constant within every tangent space at a
poi

…
\end{enumerate}
\end{proof}




\begin{exercise}
\label{exercise-kahler-form-conformal-class}
Let $X$ be a connected complex manifold of
dimension $n>1$
and let $g$ be a Kähler metric.
Show that $g$ is the only Kähler metric in
its conformal class,
i.e. if $g'=e^fg$ is Kähler then $f$ is
constant.
\end{exercise}

\begin{exercise}
\label{exercise-lefsc-injec-1-forms}
Let $X$ be a connected complex manifold and
$h$ be a Kähler metric.
Let $\omega$ be the associated Kähler form.
Show that if $\dim X \geq  2$, the wedge
product with $\omega$
is injective on 1-forms.

Show that if $\dim X \geq  2$ and $\phi$ is a
differentiable
function with values in $\mathbb{R}^+$ such
that $\phi h$ is also
a Kähler metric, then $\phi$ is constant.
\end{exercise}

\begin{proof}
Take local symplectic coordinates, i.e. such
that
$\omega = \sum_idz^i \wedge d \bar{z}^i$.
Let $\alpha$ be any 1-form and suppose it is
written
in those coordinates as
$\alpha=\sum_ia_idz^i +
\sum_i\bar{a}_id\bar{z}^i$.
Then
\begin{align*}
\alpha\wedge\omega&=
\left(\sum_ia_idz^i\right)
\wedge
\left(\sum_i dz^i \wedge d\bar{z}^i\right)
+
\left(\sum_i \bar{a}_i dz^i\right)
\wedge
\left(\sum_i dz^i \wedge d\bar{z}^i\right)\\
&=
\sum_{i,j}a_i dz^i \wedge dz^j \wedge
d\bar{z}^j
+
\sum_{i,j}\bar{a}_i d\bar{z}^i \wedge dz^j
\wedge d\bar{z}^j.
\end{align*}

\noindent
We see that if $\alpha \wedge \omega=0$
then both its $1,2$ and $2,1$ parts are zero,
which foces $a_i$ and $\bar{a}_i$ to vanish,
i.e. $\alpha=0$.
\end{proof}
\begin{exercise}
\label{exercise-holom-vecto-bundl-kahle}
\begin{reference}
\cite[Exercise 2, Chapter 3]{voi}
\end{reference}
Let $E$ be a holomorphic vector bundle of
rank $r$
over a complex manifold $X$.
Show that if $L$ is a holomorphic line bundle
on $X$,
then $\mathbb{P}(E^* \otimes L^*)=
\mathbb{P}(E^*)$
but the line bundle
$\mathcal{O}_{\mathbb{P}(E^*\otimes L^*)}(1)$
is isomorphic to
$\mathcal{O}_{\mathbb{P}(E^*)}(1) \otimes
\pi^*L$,
where $\pi:\mathbb{P}(E^*)\to X$ is the
structural morphism.
\end{exercise}

\begin{exercise}
\label{exercise-kahler-iff-i-parallel}
$(M,I,g)$ is Kähler if and only if
$\nabla^{LC}I=0$.
\end{exercise}





\section{Hodge theory}
\label{section-hodge-theory}

\noindent
On a compact hermitian manifold $(X,g)$ we
may define
$$
(\alpha,\beta):=\int_Xg_{\mathbb{C}}(\alpha,
\beta)*1,
$$
which is an hermitian product on
$\mathcal{A}_{\mathbb{C}}^*(X)$.

\begin{lemma}
\label{lemma-adjoints}
\begin{reference}
\cite[Lemma 3.2.3]{huc}
\end{reference}
Let $X$ be a compact hermitian manifold.
Then $\partial^*$ and $\bar\partial^*$
are adjoints to $\partial$ and $\bar\partial$
with respect to
$(\cdot,\cdot)$.
\end{lemma}

\begin{proof}
We have
\begin{align*}
(\partial\alpha,\beta)
&=\int_X\left<\partial\alpha,
\beta\right>_\mathbb{C}\text{Vol}
=\int_X \partial\alpha\wedge*\beta,\\
(\alpha,\partial^*\beta)
&=\int_X\left<\alpha,-
*\partial*\beta\right>_{\mathbb{C}}\text{Vol}
=\int_X\alpha\wedge\partial*\beta.
\end{align*}

\noindent
Now we also have
$\partial(\alpha\wedge\beta)
=\partial\alpha\wedge\beta+
\alpha\wedge\partial\beta$,
and the integral of the first vanishes by
Stokes theorem.
\end{proof}

\begin{lemma}
\label{lemma-harmonic-iff-closed}
Let $X$ be a compact hermitian manifold.
A form is $\bar\partial$-harmonic if and only
if
it is $\bar\partial$-closed and
$\bar\partial^*$ closed.
The same holds for $\partial$,
and in fact, on any compact oriented
Riemannian manifold $(M,g)$,
a similar result holds for the usual
$d$-Laplacian.
\end{lemma}

\begin{proof}
\begin{align*}
(\Delta \alpha,\alpha)&=((d
d^*+d^*d)\alpha,\alpha)\\
&=(d d^*\alpha,\alpha)+(d^*d\alpha,\alpha)\\
&=(d^*\alpha,d^*\alpha)+(d\alpha,d\alpha)\\
&=\|d^*\alpha\|^2+\|d\alpha\|^2.
\end{align*}
\end{proof}

\noindent
The following is the first of the symmetries
of the Hodge diamond:

\begin{lemma}
\label{lemma-hodge-commutes-with-laplacian}
\begin{reference}
\cite[Lemma A.0.14]{huc}
\end{reference}
For any Riemannian manifold $(M,g)$,
not necessarily compact, $*\Delta=\Delta*$
and $*:\mathcal{H}^k(M,g) \cong
\mathcal{H}^{m-k}(M,g)$.
\end{lemma}

\noindent
In the hermitian case, we have for a compact
hermitian connected
manifold $(M,g)$, that the pairing
\begin{align*}
\mathcal{H}^{p,q}_{\bar\partial}
\times\mathcal{H}^{n-p,n-q}_{\bar\partial}
 &\longrightarrow \mathbb{C} \\
(\alpha,\beta) &\longmapsto
\int_X\alpha\wedge\beta
\end{align*}

\noindent
is non-degenerate and thus
$$
\mathcal{H}^{p,q}_{\bar\partial}(X,g)
\cong
\mathcal{H}^{n-p,n-q}_{\bar\partial}(X,g),
$$
which is called {\it Serre duality},
the second of the symmetries of the Hodge
diamond.
We shall later generalize this Serre duality
to other vector bundles.

There is a third symmetry of the Hodge
diamond,
given by the Lefschetz operator.
Since $[L,\Delta]=0$, one of the Kähler
identities,
we see that $L$ maps harmonic forms to
harmonic forms.
Further, since $L^{n-k}:\Lambda^k(V^\vee) \to
\Lambda^{2n-k}(V^\vee)$
is bijective we have that the induced map
$$
L^{n-k}:\mathcal{H}^{p,k-p}(X,g)
\to \mathcal{H}^{n+p-k,n-p}(X,g)
$$
is injective. It is also bijective since…
thus it is an isomorphism.

\medskip\noindent
Now we are almost ready for Hodge theorem.

\begin{proposition}
\label{proposition-decom-ortho}
\begin{reference}
\cite[Proposition 3.2.2]{huc}
\end{reference}
Let $(X,g)$ be a compact hermitian manifold.
The degree, bidegree and Lefschetz
decomspositions
are orthogonal with respect to
$(\cdot,\cdot)$.
\end{proposition}

\noindent
The bidegree decomposition passes to
$\partial$ and $\bar\partial$-harmonic forms
on any compact hermitian manifold,
but only for Kähler manifolds does it pass to
$d$-harmonic forms.

\begin{proposition}
\label{proposition-decom-harmo-forms}
Let $(X,g)$ be a hermitian manifold, not
necessarily compact.
\begin{enumerate}
\item $\mathcal{H}^{k}_{\bar\partial}(X,g)
=\bigoplus_{p+q=k}\mathcal{H}^{p,q}
_{\bar\partial}(X,g)$
and
$\mathcal{H}^{k}_\partial(X,g)
=\bigoplus_{p+q=k}\mathcal{H}^{p,q}
_\partial(X,g)$.

\item If $(X,g)$ is Kähler then both
decompositions coincide with
$\mathcal{H}^{k}(X,g)
=\bigoplus_{p+q=k}\mathcal{H}^{p,q}(X,g)$.
In particular,
$\mathcal{H}^k(X,g)_{\mathbb{C}}
=\mathcal{H}^k_{\bar\partial}(X,g)
=\mathcal{H}^k_{\partial}(X,g)$.
\end{enumerate}
\end{proposition}

\noindent
We finally arrive at the Hodge decomposition.
The first result concerns only compact
Riemannian manifolds,
and ``requires some hard, but by now
standard, analysis'':

\begin{theorem}[Hodge decomposition]
\label{theorem-hodge-decomposition}
Let $(M,g)$ be a compact Riemannian manifold.
Then
$$
\mathcal{A}^k(M) =
d(\mathcal{A}^{k-1}(M))
\oplus
\mathcal{H}^k(M,g)
\oplus
d^* (\mathcal{A}^{k+1}(M))
$$
which is orthogonal with respect to
$(\cdot,\cdot)$.
Moreover, the space $\mathcal{H}^k(M)$ is
finite-dimensional.
\end{theorem}

\noindent
And now we also have
\begin{theorem}[Hodge decomposition]
\label{theorem-hodge-decomposition-hermitian}
Let $(X,g)$ be a compact hermitian manifold.
Then
$$
\mathcal{A}^{p,q}(X)=
\partial\mathcal{A}^{p-1,q}(X)
\oplus
\mathcal{H}^{p,q}_\partial(X,g)
\oplus
\partial^*\mathcal{A}^{p+1,q}(X)
$$
and
$$
\mathcal{A}^{p,q}(X)=
\bar\partial\mathcal{A}^{p-1,q}(X)
\oplus
\mathcal{H}^{p,q}_{\bar\partial}(X,g)
\oplus
\bar\partial^*\mathcal{A}^{p+1,q}(X).
$$
If $(X,g)$ is Kähler, then
$\mathcal{H}^{p,q}_\partial(X,g)=\mathcal{H}
^{p,q}_{\bar\partial}(X,g)$.
\end{theorem}

\noindent
Upshot: if $X$ is Kähler we can pass from
usual
harmonic $(p,q)$-forms to
$\bar\partial$-harmonic $(p,q)$-forms,
which takes us to $H^{p,q}(X)$.
To pass from $H^k(X)$ to $\mathcal{H}^k(X)$
and from
$\mathcal{H}^{p,q}_{\partial}(X)$ to
$H^{p,q}(X)$
we have the following two lemmas
which in fact are corollaries from the
corresponding decompositions
(Theorems \ref{theorem-hodge-decomposition}
and
\ref{theorem-hodge-decomposition-hermitian},
respectively).

\begin{lemma}
\label{lemma-de-rham-to-harmonic}
\begin{reference}
\cite[Corollary A.0.17]{huc}
\end{reference}
Let $(M,g)$ be a compact oriented Riemannian
manifold.
The natural map $\mathcal{H}^k(M) \to
H^k(M,\mathbb{R})$
that associates to any harmonic form its
cohomology class is
an isomorphism.
\end{lemma}

\begin{lemma}
\label{lemma-harmo-pq-dolbe-pq}
Let $(X,g)$ be a compact hermitian manifold.
Then the canonical projection
$\mathcal{H}_{\bar\partial}^{p,q}(X,g)\to
H^{p,q}(X)$
is an isomorphism.
\end{lemma}

\medskip\noindent
Before we pass to Hodge theorem we state one
final key result

\begin{lemma}[ddbar Lemma]
\label{lemma-ddbar}
Let $X$ be a compact Kähler manifold.
If $\alpha$ is a $d$-closed $(p,q)$-form,
it is equivalent that $\alpha$ is
\begin{enumerate}
\item $d$-exact,
\item $\partial$-exact,
\item $\bar\partial$-exact,
\item $\partial \bar\partial$-exact.
\end{enumerate}
\end{lemma}

\noindent
There is another version in terms of
$dd^c:=
I^{-1}dI=\frac{1}{2}(\partial-\bar\partial)$.

\begin{lemma}[ddc Lemma]
\label{lemma-ddc}
\begin{reference}
\cite[Lemma 3.A.22]{huc}
\end{reference}
Let $X$ be a compact Kähler manifold.
Any $d$-closed and $d^c$-exact form form is
$dd^c$-exact.
\end{lemma}



\noindent
That's all the ingredients for a
(very rough) idea on why the Hodge theorem
holds.
Upshot:
$$
H^k(X,\mathbb{C})
= \mathcal{H}^k(X)
= \bigoplus_{p+q=k}\mathcal{H}^{p,q}(X)
=
\bigoplus_{p+q=k}\mathcal{H}^{p,q}
_{\bar\partial}(X)
=\bigoplus_{p+q=k}H^{p,q}(X).
$$
\begin{enumerate}
\item Start with singular complex cohomology,
\item pass to harmonic forms
via Lemma \ref{lemma-de-rham-to-harmonic}
(this is where heavy real analysis enters;
what Huybrechts calls Hodge decomposition),
\item decompose harmonic forms in
$(p,q)$-types
via Proposition
\ref{proposition-decom-harmo-forms}
(using Kählerness),
\item use that $\mathcal{H}^{p,q}\cong
\mathcal{H}^{p,q}_{\bar\partial}$
by the same Proposition
\ref{proposition-decom-harmo-forms}
(using Kählerness),
\item go back to Dolbeault cohomology
via Lemma \ref{lemma-harmo-pq-dolbe-pq}
(again Hodge decomposition, now for
$\bar\partial$).
\end{enumerate}

\begin{theorem}[Hodge]
\label{theorem-hodge-theorem}
Let $(X,g)$ be a compact Kähler manifold.
Then there exists a decomposition
$$
H^k(X,\mathbb{C})
= \bigoplus_{p+q=k}H^{p,q}(X)
$$
which does not depend on the Kähler
structure.
Moreover, with respect to complex conjugation
on
$H^\bullet(X,\mathbb{C})=H^\bullet(X,
\mathbb{R})\otimes\mathbb{C}$
we have
$\overline{H^{p,q}(X)}=H^{q,p}(X)$,
and, by Serre duality on harmonic forms,
$H^{p,q}(X) \cong H^{n-p,n-q}(X)^\vee$.
\end{theorem}

\begin{exercise}
\label{exercise-pic-of-pn-is-z}
\begin{reference}
\cite[Exercise 3.2.11]{huc}
\end{reference}
Show that $\text{Pic}(\mathbb{P}^n) \simeq
\mathbb{Z}$.
\end{exercise}

\noindent
The next step is to check
that most results we've developed so far for
complex vector spaces and complex manifolds
descend to cohomology in Kähler manifolds.


\medskip\noindent
Let's put here the part on the cones until
I find a better place.

Recall that $\text{NS}(X)$ is the quotient
group
$\text{Pic}(X)/H^1(\mathcal{O}_X)\subset
H^2(X,\mathbb{Z})
\subset H^2(X,\mathbb{R})$
coming from the exponential exact sequence.
This has an intersection form.
In the case of a K3 we readily see that
$h^1=0$ so $\text{NS}=\text{Pic}$.
In fact, the intersection form coincides with
the usual cup product on 2-cocycles of
singular cohomology.

For a Kähler manifold we have
\begin{itemize}
\item The {\it Kähler cone}. Classes
in $H^{1,1}(X,\mathbb{R})$ that are
represented by a
Kähler form.

\item The {\it positive cone}. Classes
$\text{NS}(X)_\mathbb{R}
:=\text{NS}(X) \otimes_\mathbb{Z}\mathbb{R}$
with $\alpha^2>0$.

\item The {\it ample cone}. Classes
$\text{NS}(X)_\mathbb{R}$
with $\alpha^2>0$ and $\alpha.C>0$ for all
algebraic curves 0.
(See
\ref{algebraic-geometry-theorem-%
Nakai-Moishezon-criterion}.)
This is true, but hard to prove, see
\cite[p. 150]{huk}.
Another description is that the ample cone
is the set of classes in
$\text{NS}(X)_\mathbb{R}$
that are finite sums $\sum a_iL_i$
for $L_i \in \text{NS}(X)$ ample
and $a_i \in \mathbb{R}_{>0}$.

\item The {\it nef cone}. Classes (in
$\text{NS}(X)?)$
with $\alpha.C=\int_C\alpha\geq 0$ for any
algebraic curve.
There is also a notion of nef for analytic
cycles
which basically says that the form is
approximated
by Kähler forms.
In \cite[Theorem 4.7(ii)]{Demailly-Paun} we
have that
a $(1,1)$-class is nef (in the analytic
sense) if and only if
$\int_Y \alpha\wedge\omega^{p-1}\geq 0$ for
every
irreducible analytic set $Y \subset X$,
$p=\dim Y$
and $\bar{\omega}$ Kähler class.

\end{itemize}

\noindent
The upshot about Demailly-Pa\u un theorem
is that Kähler classes are numerically
positive
(recall that ``numerically'' means w.r.t.
intersection form)
on {\it analytic} cycles. This distinguishes
complex analytic geometry from algebraic
geometry.
This is why Demailly-Pa\u un say that
their theorem is a generalization of
Nakai-Moishezon:
because the latter is for {\it algebraic}
cycles,
indeed, found in Hartshorne's book.

The proof goes roughly as follows.
It's clear that $\mathcal{K} \subset
\mathcal{P}$
($\mathcal{P}$ is the analytically
numerically positive cone)
because, since the Kähler form
comes from the Riemannian metric
and the curves are complex,
which means that they are invariant under
the complex structure,
we have integrals
$\int_Y\omega$ for curves $Y \subset X$
but $\omega(v,Jv)>0$ for a basis $v,Jv$
of $T_pY$.
That is to say, there are no complex
Lagrangian curves.

But that's elementary.
The point of the theorem is proving that
a form in $\overline{\mathcal{K}}\cap
\mathcal{P}$,
i.e. in the closure of the Kähler cone
and which is also numerically analytically
positive,
is in fact Kähler.

So, the theorem says that the boundary of the
Kähler cone does
not intersect the numerically analytically
positive cone.

\begin{theorem}[Demailly-Pa\u un]
\label{theorem-demailly-paun}
Let $X$ be a compact Kähler manifold.
The Kähler cone is a connected component
of the cone of 1,1-classes numerically
positive on
analytic cycles.
\end{theorem}

\begin{exercise}
\label{exercise-demailly-paun}
Let $M$ be a K3 surface such that
$\text{rk}\text{NS}(M)<20$.
Using Demailly-Paun theorem,
prove that there exists a non-zero vector
$\eta \in H^{1,1}(M)$ which satisfies
$\eta^2=0$
and belongs to the boundary of the Kähler
cone.
\end{exercise}

\begin{proof}
Let $\alpha \in
\text{NS}(X)^{\perp}\setminus\{0\}$,
which is nonempty since $\dim
\text{NS}(X)<20=h^{1,1}$.
By the Hodge index theorem
\ref{theorem-hodge-index}
the intersection form on
$H^{1,1}(X,\mathbb{R})$
has index $(1,19)$.
If we assume that $X$ is projective
then we know that the restriction of
the Fubini-Study form $\omega$ is integral.
Since $\omega$ is positive (by being a Kähler
form),
we see by the index of the intersection form
that $\alpha$ must be negative.

Define $\eta_t:=t\alpha+(1-t)\omega$ for
$t\in[0,1]$.
Then $\eta_t^2=t^2\alpha^2+(1-t)^2\omega^2$.
Since at $t=0$ we have a positive number
and at  $t=1$ a negative number,
we know that there exists  $t_0$
such that $\eta_{t_0}^2=0$
by continuity of the intersection form.

Now we need to show that $\eta_{t_0}$ is
in the boundary of the Kähler cone.
For this it's enough to show that for every
$\varepsilon>0$
$\eta_{t_0+\varepsilon}$ is in the Kähler,
which
by Demailly-P\u aun
means that we have to show
that it is positive on analytic cycles.

By our choice of $\omega$ and $\alpha$
we see that
$$
\int_Y \eta_{t_0+\varepsilon}
=\int_Y
(t_0+\varepsilon)\alpha+(1-(t_0+\varepsilon))
\omega
$$
and the integral of $\alpha$ vanishes
since $\alpha$ is in the orthogonal
complement
of $\text{NS}(X)$, which contains
all analytic curves by $X$ being projective
(Chow theorem).
\end{proof}






\section{Lefschetz theorems}
\label{section-lefschetz-theorems}

\begin{theorem}[Lefschetz theorem on
(1,1)-classes]
\label{theorem-11-classes}
All the classes in $H^{1,1}(X,\mathbb{Z})
:=H^{1,1}(X) \cap H^2(X,\mathbb{Z})$ are
algebraic.
\end{theorem}

\noindent
This implies that
$\text{NS}(X):=\text{Pic}(X)/\ker c_1$
is isomorphic to $H^{1,1}(X,\mathbb{Z})$.
The ``realification'' of the integer classes,
$H^{1,1}(X,\mathbb{Z}) \otimes \mathbb{R}
=\text{NS}(X) \otimes \mathbb{R}$
is a vector subspace
$H^{1,1}(X) \cap H^2(X,\mathbb{R})$.

\begin{theorem}[Hodge index]
\label{theorem-hodge-index}
The index of the intersection pairing on
$H^{1,1}(X,\mathbb{R})$
is $(1,\rho(X)-1)$.
\end{theorem}

\begin{theorem}[Hard Lefschetz Theorem]
\label{theorem-hard-lefschetz}
\begin{reference}
\cite[p. 122]{gri}
\end{reference}
The map
$$
L^k:H^{n-k}(M) \to H^{n+k}(M)
$$
is an isomorphism.
If we define the {\it primitive} cohomology
as
$$
P^{n-k}(M) = Ker L^{k+1}
$$
we have
$$
H^m(M)=\bigoplus_k L^kP^{m-2k}(M),
$$
called the Lefschetz decomposition.
\end{theorem}

\section{Vector bundles, connections,
curvature and Chern classes}
\label{section-vector-bundles}

\begin{lemma}[Vector bundle construction from
fibers]
\label{lemma-vector-bundle-from-fibers}
\begin{reference}
\cite[Lemma 3.4]{lec}
\end{reference}
Suppose $M$ is a complex manifold,
and for each $p \in M$ we are given a
$k$-dimensional vector space $E_p$.
Let $E=\coprod_{p \in M}E_p$ and let $\pi:E
\to M$ be the obvious projection.
Suppose further that we are given
\begin{enumerate}
\item an indexed open cover $U_i$ of $M$
\item for each $i$, a bijection $
\psi_i:\pi^{-1}(U_i) \to U_i
\times\mathbb{C}^k$,
\item for each $i,j$ with $U_i\cap U_j \neq
\emptyset$
a holomorphic map $g_{ij}:U_i \cap U_J \to
\text{GL}(k,\mathbb{C})$
such that $\psi_i\circ
\psi_j^{-1}(p,v)=(p,g_{ij}(p)v$.
\end{enumerate}
\noindent
Then $\pi:E \to M$ has a unique structure as
a holomorphic vector bundle
with $(U_i,\psi_i)$ as holomorphic local
trivializations,
and with $g_{ij}$ as transition functions.
\end{lemma}

\begin{proposition}[Vector bundle
construction from cocycle]
\label{proposition-vecto-bundl-cocyc}
Let $M$ be a complex manifold and let $U_i$
be an open cover of $M$.
Suppose that for each $i,j$ such that $U_i
\cap U_j \neq \emptyset$
we are given a holomorphic map $g_{ij}:U_i
\cap U_j \to \text{GL}(k,\mathbb{C})$
such that
$$
g_{ij}(p)g_{jk}(p)g_{ki}(p)=\text{Id}
$$
for any $p \in U_i\cap U_j\cap U_k$.

Then there is a rank-$k$ holomorphic vector
bundle $E \to M$
that admits a holomorphic trivialization over
each set $U_i$
with transition functions given by $g_{ij}$.
\end{proposition}

\begin{proposition}[Holomorphic Vector Bundle
Isomorphism Criterion]
\label{proposition-isomo-vecto-bundl}
\begin{reference}
\cite[Proposition 3.7]{lec}
\end{reference}
Two vector bundles $E$ and $F$
with transition functions $g_{ij}$ and
$h_{ij}$
are isomorphic if and only if
$g_{ij}=A^{-1}_ih_{ij}A_j$
for some $A_i:U_i \to
\text{GL}_n(\mathbb{C})$.
\end{proposition}

\begin{proof}
The idea is to define a map in local
trivializations
and confirm that this glues to a global
bundle map.
This means: invent
$\varphi_i:E(U_i):=\Gamma_E(U_i) \to F(U_i)$,
i.e. a local map of sections
(yes, a map of sections, since a bundle map
is constructed as map of the total spaces,
which are points and vectors,
but we need to preserve the projection,
so all we care is about the vectors).
Then check that on $U_i\cap U_j:= U_{ij}$
we have $\varphi_i=\varphi_j$.
This literally creates a bundle map.

Try to do it for line bundles,
so $A_i$ are just functions.
Take a local frame $e_i$ for $E$
and a local frame $f_j$ for $F$.
Change coordinates and apply the map,
apply the map and change coordinates
- this should give both sides of the cocycle
condition
of the proposition.
\end{proof}

\begin{exercise}[Tensor product bundle]
\label{exercise-tensor-product-bundle}
\begin{reference}
\cite[Exercise 2.2.1]{huc}
\end{reference}
Let $E$ and $F$ be vector bundles determined
by the cocycles
$g_{ij}: U_i \cap U_j \to
\text{GL}(r,\mathbb{C})$
and $\tilde{g}_{ij}:U_i \cap U_j \to
\text{GL}(\tilde{r},\mathbb{C})$.
Verify the cocycle description for
\begin{itemize}
\item The direct sum $E \oplus F$,
\item the tensor product $E \otimes F$ and
\item the dual bundle $E^\vee$.
\end{itemize}
\end{exercise}

\begin{proof}
Let's do the tensor product.
I think that, slightly more generally than in
the
formulation of the question, if $U_i$
is a trivializing cover for $E$ and
$\tilde{U}_j$ is a trivializing cover for
$F$,
we should consider the cover $U_i\cap
\tilde{U}_j$.
Let us denote such a cover simply by $U_i$.

Define a cocycle by
\begin{align*}
h_{ij}: U_i\cap U_j &\longrightarrow
\text{GL}(r\tilde{r},\mathbb{C}) \\
p &\longmapsto g_{ij}(p)\otimes
\tilde{g}_{ij}(p).
\end{align*}

\noindent
The key observation is that while formally
the image of $h_{ij}$
lies in $\text{GL}(\mathbb{C}^r \otimes
\mathbb{C}^{\tilde{r}})$,
we can regard such a space as
$\text{GL}(r\tilde{r},\mathbb{C})$
by defining, for $A \in
\text{GL}(r,\mathbb{C})$
and $A' \in \text{GL}(\tilde{r},\mathbb{C})$,
the map $A(v^ie_i) \otimes
\tilde{A}(\tilde{v}^j\tilde{e}_j)
=v^i\tilde{v}^j Ae_i \otimes
\tilde{A}\tilde{e}_j$
where $e_i$ and $\tilde{e}_j$ are bases
of $\mathbb{C}^r$ and
$\mathbb{C}^{\tilde{r}}$.
This is indeed a linear transformation of the
$r\tilde{r}$-dimensional $\mathbb{C}$-vector
space
generated by the basis $e_i \otimes
\tilde{e}_j$.

Thus the cocycle condition follows simply by
the cocycles of $E$ and $F$, and
commutativity of the
coordinate functions.
\end{proof}

\begin{example}[Determinant bundle]
\label{example-determinant-bundle}
Let $\pi:E \to M$ be a vector bundle of rank
$r$.
The {\it determinant bundle} of $E$ is given
by the following data.
A space $\tilde{E}:=\coprod_{p \in
M}\Lambda^rE_p$
with the obvious projection map
and topology induced from such map.
Since we have a cover $U_i$ of $M$
along with local trivializations
$U_i\overset{\varphi_i}{\simeq}
U_i\times\mathbb{C}^r$
we may just postcompose $\varphi_i$ with the
functor $\Lambda^r$
on the second entry, giving trivializations
$$
\xymatrix{
\pi(U_i)\ar[r]^-{\varphi_i}_-{\simeq}
& U_i \times
\mathbb{C}^r\ar[r]^-{\text{Id}\times
\Lambda^r}
& U_i \times \Lambda^r\mathbb{C}^r=U_i \times
\mathbb{C}
}
$$
We can also apply $\Lambda^r$ to the
transition functions
$g_{ij}$ on the intersections
$U_{ij}:=U_i\cap U_j$
$$
\xymatrix{
U_{ij}\ar[r]^-{\varphi_i}&U_{ij} \times
\mathbb{C}^r
\ar[d]^{g_{ij}\in\text{GL}_r(\mathbb{C})}\\
U_{ij}\ar[r]^-{\varphi_j}&U_{ij}
\times\mathbb{C}^r
}
$$
giving maps in $\text{GL}_1(\mathbb{C})$.
Recall that $\Lambda^r$ is functorial on
linear maps
$g_{ij}\in \text{GL}_n(\mathbb{C})$ because
it is given by the determinant of $g$,
and the determinant is multiplicative.
We thus obtain the cocycle condition on the
maps $g_{ij}$
after applying $\Lambda^r$.
\end{example}

\begin{definition}
\label{definition-connection-vector-bundle}
Let $E$ be a vector bundle.
A {\it connection} on $E$ is an
$\mathbb{R}$-linear map
$\nabla:H^0(E)\to H^0(E) \otimes \Omega^1(M)$
satisfying the identity
$\nabla(fs)=dfs + f\nabla s$.
\end{definition}

\begin{definition}
\label{definition-hermitian-connection}
\begin{reference}
\cite[Definition 4.2.9]{huc}
\end{reference}
Let $(E,h)$ be a hermitian vector bundle.
A connection $\nabla$ on $E$ is called
{\it hermitian} if
$d(h(s,s'))=h(\nabla s,s')+h(s,\nabla s')$.
\end{definition}

\noindent
Perhaps this notion defined on a seminar
talk by Beatrice Brienza is the same:
let $(M,I,g)$ be a hermitian manifold. A
linear connection $\nabla$ on $TM$ is
{\it hermitian} if
$\nabla g=0$ and $\nabla I =0$.

\begin{definition}
\label{definition-holomorphic-connection}
A connection $\nabla$ on a vector bundle $E$
is said to be {\it compatible with the
complex structure $I$} if
$\pi^{0,1}\nabla s=\bar\partial s$
where $\pi^{0,1}:\mathcal{A}^1(M)
\to \mathcal{A}^{0,1}(M)$ is the projection.
\end{definition}

\begin{proposition}
\label{proposition-chern-connection}
On a holomorphic vector bundle
equipped with a hermitian metric there exists
a unique hermitian connection
compatible with the complex structure
called the {\it Chern connection}.
\end{proposition}

\medskip\noindent
Now we discuss line bundles on
$\mathbb{P}^n$.

\begin{definition}
\label{definition-line-bundles-in-pn}
$$
\mathcal{O}(1):=\mathcal{O}(-1)^\vee,\qquad
\mathcal{O}(d):=\bigotimes_d \mathcal{O}(1).
$$
\end{definition}

\noindent
To show that the sections of $\mathcal{O}(d)$
are degree-$d$ homogeneous polynomials we
first
show that the fibers are degree-$d$
homogeneous functions
(not necessarily polynomials)
and then pass to the global situation.

\begin{lemma}
\label{lemma-fiber-od-homog-degre-d-funct}
The fiber $\mathcal{O}(d)_p$ is in bijection
with
degree-$d$ homogeneous functions on $p$.
\end{lemma}

\begin{proof}
So far we can only say that the fibers of
$\mathcal{O}(-1)\otimes \mathcal{O}(1)$ are
the fibes of $\mathcal{O}$.
This means that for a fixed $f \in
\mathcal{O}(1)_p$
we have a linear functional $z \mapsto
f\otimes z$
on $\mathcal{O}(-1)_p$. Any linear functional
on $\mathcal{O}(-1)_p$
can be realised this way.

To generalize this for degree $d$ we do as
follows.
First note that
$$
\Hom(\mathcal{O}(-1)_p,\mathbb{C})
\otimes\Hom(\mathcal{O}(-1)_p,\mathbb{C})
\cong
\Hom(\mathcal{O}(-1)_p\otimes\mathcal{O}(-1)
_p,\mathbb{C}).
$$
(This is not true for general modules,
but in this case it is using a basis.)

By the universal property of tensor product
(see Commutative Algebra Theorem
\cite{theorem-unive-prope-tenso-produ})
we know that
$\Hom(\mathcal{O}(-1)_p \otimes … \otimes
\mathcal{O}(-1)_p,\mathbb{C})$
is the set of multilinear maps
$\text{Lin}(\mathcal{O}(-1)_p
\times…\times\mathcal{O}(-1)_p,\mathbb{C})$.
Any such map $L$ gives a degree-$d$ function
(a polynomial) $P$ by
$P(x):=L(x,…,x)$.
In fact, since the fiber $\mathcal{O}(-1)_p$
is one-dimensional,
we can reconstruct $L$ from $P$ by defining
the multilinear map on a basis (which is
composed by a single vector $x$)
by $L(x,…,x):=P(x)$. This independs on the
choice
of basis since if $x'=\lambda x$ and we
define $L'(x',…,x'):=P(x')$,
We see that $\lambda^d L'(x,…,x)=L'(x',…,x')=
=P(x')=P(\lambda x)=\lambda^d P(x)=\lambda^d
L(x,…,x)$.
\end{proof}

\noindent
To pass from the local situation in Lemma
\ref{lemma-fibers-of-od}
to a section we do as follows.

\begin{exercise}
\label{exercise-polynomials-are-sections}
\begin{reference}
\cite[Exercise 2.2.8]{huc}, and
\cite[Lemma 3.30]{lec} for the answer.
\end{reference}
Show that any non-trivial homogeneous
degree-$d$ polynomial $0 \neq f \in
\mathbb{C}[z_0,…,z_n]_d$ can be considered as
a non-trivial section of $\mathcal{O}(d)$ on
$\mathbb{P}^n$.
\end{exercise}

\begin{proof}
Let $P$ be a homogeneous degree-$d$
polynomial. We define a section fiberwise and
then show it is holomorphic. For $p \in
\mathbb{P}^n$ we restrict $P$ to the
1-dimensional vector space that $p$ is in
$\mathbb{C}^{n+1}$. To show this is
holomorphic we can evaluate our section on a
local frame for $\mathcal{O}(d)$, which gives
a holomorphic function since $P$ is a
polynoimal (see \cite[Proposition 3.35]{lec}.
\end{proof}

\begin{theorem}
\label{theorem-secti-od-degre-d-homog-poly}
Sections of $\mathcal{O}(d)$ are in
correspondence with homogeneous degree-$d$
polynomials.
\end{theorem}

\begin{proof}
It only remains to prove the direct
implication.
Let $s \in H^0(\mathcal{O}(d))$.
Then $s(p) \in \mathcal{O}(d)_p$ is a
homogeneous degree-$d$ function $f_p$
by Lemma
\ref{lemma-fiber-od-homog-degre-d-funct}
defined on the vector space $p$.

Now define a holomorphic homogeneous
degree-$d$ function $f$ on
$\mathbb{C}^{n+1}\setminus\{0\}$ as follows.
For any $z \in
\mathbb{C}^{n+1}\setminus\{0\}$ we can take
$p$
to be the line passing through $z$
and assign the number $f_p(z)$
(since $z \in p$).

Finally we show that $f$ must be a polynomial
by considerations on its Taylor series
after extending to $\mathbb{C}^{n+1}$ via
Hartog's theorem
(see \cite[Theorem 3.36]{lec}).
\end{proof}


\begin{exercise}
\label{exercise-line-bundle-trivial-iff}
\begin{reference}
\cite[Exercise 2.2.5]{huc}
\end{reference}
Let $L$ be a holomorphic line bundle on a
compact complex manifold $X$.
Show that $L$ is trivial if and only if $L$
and its dual $L^\vee$
admit non-trivial global sections.
\end{exercise}

\begin{proof}
Direct implication is easy: trivial bundle
has 1-dimensional space of sections.
For the converse choose nonvtrivial sections
$s$ of $L$
and $s^\vee$ of $L^\vee$.
Then $s \otimes s^\vee$ which is a section
of the trivial bundle.
If such a section has a zero, then it must be
zero throughout $X$.
\end{proof}

\medskip\noindent
Now let's discuss Serre's duality.
It is a consequence of Hodge theorem,
which says that Hodge star operator descends
to cohomology.
In fact, the general result is

\begin{theorem}[Hodge]
\label{theorem-hodge-star-descends}$$
*: H^{p,q}(X) \longrightarrow H^{n-p,n-q}(X)
$$
is an isomorphism.
\end{theorem}

\noindent
However, we shall not use this isomorphism
directly for Serre duality, see \cite[Remark
4.1.17]{huc}. What we do instead is show that
the pairing of forms given by integration is
nondegenerate. The upshot is that the proof
of nondegeneracy is uses Hodge star
operator in cohomology.

\begin{proposition}[Serre duality]
\label{proposition-serre-duality}
\begin{reference}
\cite[Proposition 4.1.15]{huc}
\end{reference}
Let $X$ be a compact complex manifold.
For any holomorphic vector bundle $E$ on $X$,
the natural pairing
$$
H^{p,q}(X,E) \times H^{n-p,n-q}(X,E) \to
\mathbb{C}
$$
is nondegenerate.
\end{proposition}

\begin{proof}
We must show that for any $\alpha \in
H^{p,q}(X,E)$ there exists $\beta \in
H^{n-p,n-q}(X,E)$ such that
$$
\int_X\alpha\wedge \beta\neq 0.
$$
Choose
$$
\beta=\overline{*}\beta
$$
where $\overline{*}$ is the conjugate of the
Hodge star operator in cohomology.
\end{proof}

\noindent
Of course, this proposition lacks context in
these notes; we should properly extend Hodge
theory to sections of arbitrary sections of
vector bundles; but it looks as if the theory
for sections of the cotangent algebra
generalizes without the need of introducing
any new ideas.

We put $p=0$ to obtain
$$
H^{0,q}(X)=H^q(X,\Omega^0)=H^q(X,\mathcal{O})
$$
and
$$
H^{n,n-q}(X)=H^{n-q}(X,\Omega^n)
=H^{n-q}(X,\omega_X)
$$
that is, we have the usual isomorphism
$$
H^q(X,\mathcal{O}_X)
\cong H^{n-q}(X,\omega_X)^\vee.
$$
Tensoring with a sheaf
$\mathcal{F}$ gives (though
I'm not sure at the time how
this machinery works…)
$$
H^q(X,\mathcal{F})
\cong
H^{n-q}(X,\mathcal{F}^\vee \otimes
\omega_X)^\vee.
$$


\section{Line bundles and divisors}
\label{section-line-bundles}

\noindent
A simple way of associating a line bundle to
an irreducible
hypersurface $Y\subset X$ is by noting that
by Weierstrass Division Theorem
\ref{complex-analysis-theorem-%
weierstrass-division-theorem}
the ideal sheaf of $Y$ is locally generated,
i.e. there is an open cover $U_i$ of $X$
such that $\mathcal{I}(U_i)=(f_i)$.
This means that at intersections we have
$f_j \in (f_i)$ and $f_i\in (f_j)$,
which says that $f_i$ and $f_j$ are
associates,
i.e. there is a unit $u_{ij} \in
\mathcal{O}^*_X(U)$
(the units in the ring of holomorphic
functions
are nowhere vanishing functions)
such that $f_i=u_{ij}f_j$.
Thus $u_{ij}:U_i\cap U_j \to
\text{GL}(1,\mathbb{C})$
define a cocycle in $H^1(X,\mathcal{O}^*_X) =
\text{Pic}(X)$.

\begin{definition}
\label{definition-divisor}
A {\it Weil divisor} is a formal sum
$$
D=\sum a_i Y_i
$$
where $a_i\in \mathbb{Z}$ and $Y_i$ are
irreducible hypersurfaces.
\end{definition}

\noindent
An irreducible hypersurface $Y$ clearly
defines a Weil divisor,
and so does a reducible reduced hypersurface
with irreducible
components $Y_i$ putting $a_i=1$.
To allow for coefficients $>1$ we may admit
nonreduced hypersurfaces,
so that we have larger multiplicites.

This can also be presented as follows.
Given a holomorphic function $g$, we define
{\it order of vanishing of $g$ at $P \in Y$}
to be the largest number $\text{ord}_{P,Y}g$
such that
$$
g=f_i^{\text{ord}_{P,Y}g}h,\qquad h \in
\mathcal{O}^*_{X,P},
$$
where $f_i$ is an irreducible function
vanishing along $Y$ near $P$,
i.e. a generator of the local ring of $Y$.
Further, by coherence/preservation relative
primes in nearby stalks
(see \cite[p. 696]{gri})
the order of $g$ is constant for any $P \in
Y$.

Thus for any holomorphic function $g$
we can define the divisor
$$
\text{div}(g):=\sum_{Y}\text{ord}_Y(g)Y
$$
where the sum ranges in all irreducible
hypersurfaces of $X$.

Similarly, for a meromorphic function
$g=g_1/g_2$
we define the order of vanishing of $g$ at
$Y$ to be
$$
\text{ord}_Y(g)=\text{ord}_Y(g_1)-\text{ord}
_Y(g_2).
$$
The associated divisor to a meromorphic
function
is defined identically as that of a
holomorphic function,
and we observe that there might appear
negative coefficients.
Divisors obtained in this way from
meromorphic functions
are called {\it principal divisors}.

This shows how to assign a divisor from any
meromorphic function,
i.e. it's a map
$$
H^0(X,\mathcal{M}^*) \to \text{Div}(X).
$$
Note that if $g$ is in fact a nowhere
vanishing meromorphic or
holomorphic function,
the corresponding divisor is trivial,
i.e. it is in the kernel of our map.
In fact,
$$
H^0(X,\mathcal{M}^*/\mathcal{O}^*) \simeq
\text{Div}(X).
$$
Indeed, given a divisor we can assign a
section of
$\mathcal{M}^*/\mathcal{O}_X^*$ as follows.
If the divisor is an irreducible hypersurface
locally given by $V(f_i)$
we assign the meromorphic function given
locally as $f_i$.
If the divisor has several irreducible
connected components
we consider the product of the irreducible
generators
at every point.

\medskip\noindent
It remains to connect $\text{Div}(X)$ and
$\text{Pic}(X)$.
The idea is that a section $f$ of
$\mathcal{M}^*/\mathcal{O}_X^*$
defines a cocycle in $H^1(\mathcal{O}_X^*)$.
Indeed, $f$ is given, by definition of
quotient sheaf (?),
by the data of an open cover $U_i$ and
meromorphic functions
$$
f_i=\frac{g_i}{h_i} \in \mathcal{M}^*
$$
which coincide at the intersections $U_i \cap
U_j$,
that is,
$$
f_i|_{U_i\cap U_j} = f_j|_{U_i\cap U_j}.
$$
But since we are in the quotient ring,
equality means that
$$
f_i f_j^{-1} \in \mathcal{O}_X^*.
$$
This defines a cocycle, and thus a map
$$
\text{Div}(X) \to \text{Pic}(X).
$$
This map can also be though of as induced
from
$$
\xymatrix{
0\ar[r]
&\mathcal{O}^*_X\ar[r]
&\mathcal{M}^*\ar[r]
&\mathcal{M}^*/\mathcal{O}_X^*\ar[r]
&0
}
$$
via the cohomology long exact sequence giving
$$
\xymatrix{
\cdots\ar[r]
& H^0(\mathcal{M}^*_X)\ar[r]
&H^0(\mathcal{M}^*_X/\mathcal{O}^*_X)\ar[r]
&H^1(\mathcal{O}_X^*)\ar[r]
&\cdots
}
$$
Indeed, the kernel of the map $\text{Div}(X)
\to \text{Pic}(X)$
is given by the principal divisors,
i.e. those coming from meromorphic sections.

In general there is no way to ensure this map
is an isomorphism.

\medskip\noindent
Now we discuss how to obtain submanifolds
from sections.
We will see that $\text{Div}(X)$
is isomorphic to the group of line bundles
equipped with a meromorphic section.


Bear in mind the following observation for a
rank-$r$
vector bundle $E \to X$.
Suppose $(\varphi_i,U_i)$ are local
trivializations of $E$,
i.e. $\varphi_i:\pi^{-1}(U_i)
\xrightarrow{\simeq}U_i \times \mathbb{C}^k$.
Take a local section $s :U_i \to
\pi^{-1}(U_i)$.
Take a local basis $e_i$ of $\mathbb{C}^k$.
Then
$$
s_i(x)=f_i^k(x)\varphi_i^{-1}(x,e_k)
$$
for some functions $f_i^k:U_i \to
\mathbb{C}$.

At an intersection $U_i \cap U_j$,
we have
$$
s_i(x)=f_i^k(x)\varphi_i^{-1}(x,e_k)
=
s_j(x)=f_j^k(x)\varphi_j(x,e_k).
$$
Applying $\varphi_i$ we see that
$$
f_i^k e_k = f_j^k g_{ij}(e_k).
$$

In the case of a line bundle this just says
that sections
are given by local functions $f_i:U_i \to
\mathbb{C}$ satisfying
$$
f_i=g_{ij}f_j
$$
where $g_{ij}$ is the cocycle defining the
line bundle.

Let $L$ be a line bundle.
A {\it meromorphic section} of $L$ is a
section
of $\mathcal{O}(L)
\otimes_\mathcal{O}\mathcal{M}$,
which is given by a collection of meromorphic
functions
$s_i \in \mathcal{M}(U_i)$
such that
$$
s_i=g_{ij}s_j.
$$
That is, sections of the line bundle viewed
as a module over $\mathcal{O}$
with coefficients that can be meromorphic
functions.

The orders of the coefficient meromorphic
functions
is well defined along hypersurfaces of $X$,
and thus
we can define the
{\it divisor $\text{div}(s)$ associated to
the section $s$ of $L$}
by
$$
\text{div}(s)
:= \sum_Y\text{ord}_Y(s)Y.
$$

Conversely, given any divisor $D$
we can assign a line bundle $L$ as before,
namely if the divisor is given by the
meromorphic functions $f_i$
we put $g_{ij}=f_i/f_j$.
But now we observe
the existence of a distinguished section $s$
such that $\text{div}(s)=L$.
Indeed, we note that
$$
f_i=f_if_j^{-1}f_j=g_{ij}f_j,
$$
which is a section of $L$,
and in fact (check…) $\text{div}(s)=L$.

We conclude that $L$ is the line bundle
associated to a divisor
if and only if it has a meromorphic section
not identically zero.
It is the line bundle associated to an
effective divisor
if and only if it has a holomorphic section
not identically zero
--- whose zero set is the divisor.

\medskip\noindent
Upshot. We have more than the submanifold
associated to a section:
we have the divisor associated to a
meromorphic section,
which gains to becoming less geometrically
intuitive
the virtue of being an element of a group.

\medskip\noindent
Now we turn to line systems.
Take a divisor
$$
D=\sum a_iV_i.
$$
The divisor
$$
D+(f)
$$
is linearly equivalent to $D$ --- their
difference is principal.
If
$$
D+(f)\geq 0,
$$
then $f$ must be holomorphic away from $V_i$,
and on $V_i$ it can have poles with the
condition
$$
\text{ord}_{V_i}(f) \geq  -a_i.
$$
Denote $\mathcal{L}(D)$ the space of such
functions.

Note that if $s_0$ is a section such that
$(s_0)=D$, then multiplication by $s_0$ gives
an identification
$$
\mathcal{L}(D) \xrightarrow{\otimes s_0}
H^0(X,L).
$$

\begin{definition}
\label{definition-line-system}
We denote by
$$
|D|\subset \text{Div}(X)
$$
the set of all effective divisors linearly
equivalent to $D$.
\end{definition}

\noindent
$\mathcal{L}(D)=H^0(X,L)$ is almost the same
thing as $|D|$.
At least in the compact case,
two functions in $f,f' \in \mathcal{L}(D)$
give the same divisor when $f=\lambda f'$
for some constant function $\lambda \in
\mathbb{C}$.

Thus
$$
|D|=\mathbb{P}H^0(X,L).
$$
This is called a {\it complete linear
system}.
More generally, a {\it linear system} is a
linear
subspace of $\mathbb{P}H^0(X,L)$.
The {\it base locus} of $|D|$
is the set of points where all the global
sections
generating $H^0(X,L)$ vanish.
We will see in Section
\ref{section-kodaira-embedding}
the relationship between line systems and
embeddings to projective space.

\begin{theorem}[Bertini]
\label{theorem-bertini}
The generic divisor of a linear system
is smooth away from the base locus.
\end{theorem}

\medskip\noindent
Now let's introduce the first adjunction
formula,
which says that when a divisor is a smooth
hypersurface,
it's associated line bundle is its
holomorphic normal bundle.

\begin{theorem}[Adjunction formula I]
\label{theorem-adjunction-formula-i}
\begin{reference}
\cite[p.146]{GH}
\end{reference}
If $X$ is a smooth hypersurface of a smooth
complex manifold $Y$, then
\begin{equation}
\label{equation-adjunction-i}
\mathcal{O}_Y(X)|_{X} \simeq
\mathcal{N}_{X/Y},
\end{equation}
where $\mathcal{N}_{X/Y}:=T_Y'/T_X'$.
\end{theorem}

\begin{proof}
If $X$ is locally defined
by the vanishing of $f_i$ on an open cover
$U_i$ of $Y$,
then $\mathcal{O}_Y(X)$ is given by the
cocycle $f_{ij}:=f_if_j$.

To prove Equation \ref{equation-adjunction-i}
it's enough to show that
$\mathcal{O}_Y(X)|_{X}\otimes
\mathcal{N}^\vee_{X/Y}$
is trivial, which can be done by constructing
a nonvanishing section.
Such a section is given by
$$
df_i=d(f_{ij}f_j)=df_{ij}\cdot
f_j+f_{ij}df_j,
$$
since the first term vanishes along $X$
(because $f_j$ does)
and the remaining terms define precisely a
section
of ${T'}^\vee_Y/{T'}^\vee_X$,
the conormal bundle ---
indeed, sections of such a bundle
are holomorphic differencials vanishing along
$X$.
\end{proof}

\noindent
Now let's find the canonical bundle of
$\mathbb{P}^n$.
By the above discussion, we need the inverse
map of
$$
\text{Div}(\mathbb{P}^n) \to \{L \in
\text{Pic}(X):
\text{$L$ has a nonzero meromorphic
section}\},
$$
that is, we need to find a nonzero
meromorphic section $s$
of the canonical bundle $\Omega^{n,0}=K$.
Then the divisor will be given by
$$
\sum_Y \text{ord}_Ys [Y].
$$
A meromorphic section is given by local
sections $s_i$
so that
$$
s_i=g_{ij}s_j
$$
where $g_{ij}$ is the cocycle of $K.$

Consider in the canonical chart $z_0\neq 0$
of $\mathbb{P}^n$
the local meromorphic section of $K$
$$
s_0=
\frac{dw_1}{w_1}\wedge\frac{dw_2}{w_2}\wedge…
\wedge\frac{dw_n}{w_n}
$$
where $w_i=z_i/z_0$ are the coordinate
functions in $z_0\neq 0$.

In any chart $z_j\neq 0$ we obtain the local
expression
$$
(-1)^j\frac{dw_0'}{w_0'}
\wedge…
\wedge \widehat{\frac{dw_j'}{w_j'}}
\wedge…
\wedge \frac{dw_n'}{w_n'}
$$
according to the canonical change of
coordinates, which says
$$
w_i=\frac{w_i'}{w'0}, i\neq j,\qquad
w_j=\frac{1}{w_0'}.
$$

The point is that we constructed a
meromorphic section which has
a single pole along each hyperplane $z_i=0$,
which means
$$
K=-(n+1)H.
$$
\medskip\noindent
The next adjunction formula is the key to
computing
the canonical bundle of submanifolds.
It is a consequence of this linear algebra
fact:

\begin{lemma}
\label{lemma-detrminant-of-ses}
If
$$
\xymatrix{
0\ar[r]
&E'\ar[r]
&E\ar[r]
&E''\ar[r]
&0
}
$$
is a short exact sequence of vector bundles
(of possibly different ranks),
then
$$
\det(E)\cong \det(E')\otimes \det(E'').
$$
\end{lemma}

\begin{proof}
First notice that each of the three
determinant bundles gives fiberwise the
top-exterior
power of the corresponding vector spaces,
thus varying the rank of the exterior power
taken.
Then define a map
\begin{align*}
x_1 \wedge…\wedge x_{\text{rk}E'}
&\otimes x_{\text{rk}E'+1}\wedge…\wedge
x_{\text{rk}E'+\text{rk}E''}\\
&\qquad \qquad \mapsto x_1 \wedge…\wedge
x_{\text{rk}E'}
\wedge x'_{\text{rk}E'+1}\wedge…\wedge
x'_{\text{rk}E'+\text{rk}E''}
\end{align*}
where $x'_{\text{rk}E'+i}$ maps to
$x_{\text{rk}E'+i}$
under the surjection given by the exact
sequence.

It's clear that such a map is injective:
if any monomial in the image is zero, then
one of the
factors must be zero, and then the element
we began with also vanishes.
Being two vector spaces of dimension 1,
the map must also be surjective.
Check this map is well-defined!
This gives an isomorphism in every fiber.

In fact the isomorphism above acts like 1.
To show that this glues to a bundle
isomorphism
we have to check some cocycle condition
as in Proposition
\ref{proposition-isomo-vecto-bundl}.
In this case the transformation functions
$A_i$
are just 1 by definition of the morphism.
The cocycle condition follows from the fact
that
any short exact sequence splits,
so on $U_{ij}$ we have the isomorphism
$E(U_{ij}) \simeq E'(U_{ij})\oplus
E''(U_{ij})$.
Recall that the cocycle of the direct sum of
bundles is given by
$\begin{pmatrix}h'_{ij}&0\\
0&h''_{ij}\end{pmatrix}$
so that when taking determinants
we obtain that
$g_{ij}=h'_{ij}h''_{ij}:=h_{ij}$
(since taking determinant bundle is taking
determinant of
the transition functions).
Thus the transition functions of the two
bundles in question
actually coincide,
and the cocycle condition in Proposition
\ref{proposition-isomo-vecto-bundl}
is the trivial equality $g_{ij}=h_{ij}$.
\end{proof}

\noindent
Here's a cleaner version of this proof:

\begin{proof}
Choose an open cover $\{U_i\}$ on which the
sequence splits, so that
\[
E|_{U_i}\simeq E'|_{U_i}\oplus E''|_{U_i}.
\]
Pick local frames $(e'_{i,1},\dots,e'_{i,r})$
for $E'|_{U_i}$ and
$(e''_{i,1},\dots,e''_{i,s})$ for
$E''|_{U_i}$, and use the splitting to obtain
a frame $(e'_{i,\bullet},e''_{i,\bullet})$ of
$E|_{U_i}$.

On overlaps $U_{ij}$ the transition matrix of
$E$ in these adapted frames has
block diagonal form
\[
g^E_{ij}=
\begin{pmatrix}
g'_{ij} & 0\\
0 & g''_{ij}
\end{pmatrix},
\]
hence
\[
\det(g^E_{ij})=\det(g'_{ij})\,\det(g''_{ij}).
\]
Thus the transition functions of $\det(E)$
coincide with those of
$\det(E')\otimes\det(E'')$, giving an
isomorphism of line bundles. In these
frames it is given by
\[
(e'_{i,1}\wedge\cdots\wedge e'_{i,r})\otimes
(e''_{i,1}\wedge\cdots\wedge e''_{i,s})
\longmapsto
e'_{i,1}\wedge\cdots\wedge e'_{i,r}\wedge
e''_{i,1}\wedge\cdots\wedge e''_{i,s}.
\]

If a different local splitting is chosen, the
adapted frames differ by a
matrix of the form
\[
\begin{pmatrix}
I & B\\
0 & I
\end{pmatrix},
\]
whose determinant is $1$. Therefore the
induced map on determinant lines is
unchanged, showing that the above isomorphism
is independent of the splitting
and hence canonical.
\end{proof}

\begin{theorem}[Adjunction formula II]
\label{theorem-adjunction-ii}
Let $X$ be a complex submanifold of a complex
manifold $Y$.
Then
\begin{equation}
\label{equation-adjunction-ii}
(K_Y\otimes \mathcal{O}_Y(X))|_{X}\cong K_X.
\end{equation}
\end{theorem}

\begin{proof}
From the holomorphic tangent bundle exact
sequence we get
$$
\xymatrix{
0\ar[r]
&T'_X\ar[r]
&T'_Y|_{X}\ar[r]
&N_{X/Y}:=T'_Y/T'_X\ar[r]
&0
}
$$
Taking determinant of this exact sequence we
obtain the result.
\end{proof}


\medskip\noindent
Now we turn to Chern classes of line bundles.

Slogan: the first Chern class of a line
bundle
is the Poincaré dual of the divisor's
fundamental class
and the cohomology class of the curvature
form
of any connection of $L$.
It is an integral $(1,1)$-curvature form.

The {\it curvature form} is the differential
2-form
that represents the curvature operator with
respect to a given connection
on the line bundle.
This can also be defined for affine vector
bundles of any rank.

\begin{proposition}
\label{proposition-first-chern-class}
\begin{reference}
\cite[p. 141]{gri}
\end{reference}
\begin{enumerate}
\item For any line bundle $L$ with curvature
form $\Theta$,
$$
c_1(L)=\left[\frac{\sqrt{-1}}{2\pi}
\Theta\right]
\in H^2_{dR}(X).
$$

\item If $L=[D]$ for some $D \in
\text{Div}(M)$,
$c_1(L)$ is the Poincaré dual of the
fundamental class of $D$.
\end{enumerate}
\end{proposition}

\medskip\noindent
Finally let's just put some more exercises
here.

\begin{exercise}[Curves on surfaces and the
Hodge index theorem]
\label{exercise-curve-surfa-hodge-index}
\begin{reference}
\cite[Exercise 2, Chapter 11]{voi}
\end{reference}
Let $S$ be a smooth projective surface and
let $\phi:S \to \mathbb{P}^n$
be a holomorphic map satisfying the following
properties:
\begin{itemize}
\item There exists a curve $C \subset S$
such that $\phi(C)$ is a finite set of
points.
(Here $C$ is a not necessarily connected
finite union
of irreducible hypersurfaces $C_i$ which are
not necessarily smooth.)

\item There exists an open set of $S$ where
$\phi$ is an embedding.
\end{itemize}
\begin{enumerate}
\item Show that the class
$h=\phi^*c_1(\mathcal{O}_{\mathbb{P}^n}(1))$
satisfies $\left<h,h\right> >0$ where
$\left<\cdot,\cdot\right>$ is the
intersection form on $H^2(S,\mathbb{Z})$.

\item Show that $\left<[C_i],h\right>=0$ for
any $i$.
Deduce from the Hodge index theorem that the
intersection form
$\left<\cdot,\cdot\right>$ is negative
definite on the subgroup
of $H^2(S,\mathbb{Z})$ generated by the
$[C_i]$'s.

\item Show that $\left<C_i,C_j\right>\geq 0$
for $i\neq j$.
Deduce then from the above that the $[C_i]$'s
are independent in $H^2(S,\mathbb{Z})$.
\end{enumerate}
\end{exercise}

\begin{proof}
\begin{enumerate}
\item Here's an outline of how to prove this
result.
First notice that $h$ is the class of a line
bundle
since $\varphi^*c_1=c_1\varphi^*$.
Thus $\left<h,h\right>$ is the intersection
number
of two divisors in the class $h$.

To compute their intersection number
we use that $\varphi$ is an embedding on an
open set of $S$,
say $U \simeq \varphi(U) \subset
\mathbb{P}^n$.
To compute the intersection number we take
two generic hyperplanes $H_1$ and $H_2$
in $\mathbb{P}^n$ and intersect them with
$\phi(U)$.
The dimension of the resulting submanifold is
zero.
Thus the intersection is a finite number of
points.

\item Since $C_i$ is mapped to a set of
points in $\mathbb{P}^n$,
we compute the intersection number as above,
which gives zero
for generic choice of $H_1,H_2$. Hodge index
theorem tells us that
the signature of $\left<\cdot,\cdot\right>$
is $(1,\rho-1)$ on
$H^2(S,\mathbb{Z})$. Since $h$ is positive
and $C_i$ is orthogonal to $h$,
it's necessary that $C_i$ are negative
vectors.

\item This is by definition/construction of
the intersection form on a surface:
for two irreducible divisors of the surface,
the intersection number is the number of
intersection points
of the divisors counted with multiplicity.
Such a number cannot be negative.
The fact that the $C_i$'s are independent
follows from a linear-algebra
argument.
\end{enumerate}
\end{proof}

\section{Kodaira Vanishing theorem}
\label{section-kodaira-vanishing}

\noindent
In \cite[p. 150]{gri}
we see that the curvature form of the
tautological bundle
is given by $\partial\bar\partial \text{log}
|z|^2$
where $z$ is a nonzero section.
Not giving further details for now,
this implies that the Chern class of the
hyperplane bundle
is the Fubini-Study form in $\mathbb{P}^n$.
We already discussed that this form is
positive
in the sense of Definition
\ref{definition-positive-11form}.

\begin{definition}
\label{definition-positive-line-bundle}
\begin{reference}
\cite[p. 148]{gri}
\end{reference}
Let $M$ be a compact Kähler manifold.
A line bundle $L \to M$ is {\it positive}
if there exists a metric on $L$ with
curvature form $\Theta$
such that $(\sqrt{-1}/2\pi)\Theta$ is a
positive $(1,1)$-form.
\end{definition}

\begin{theorem}[Kodaira-Nakano Vanishing
Theorem]
\label{theorem-kodaira-nakano-vanishing}
\begin{reference}
\cite[p. 154]{gri}
\end{reference}
[Let $M$ be a compact Kähler manifold of
complex dimension $n$.]
If $L \to M$ is a positive line bundle,
then
$$
H^q(M,\Omega^p(L)) = 0,\qquad
p+q>n=\dim_\mathbb{C} M.
$$
\end{theorem}

\noindent
Here's the same result as it appears in
\cite[sea], where it is not proved.

\begin{theorem}[Kodaira Vanishing]
\label{theorem-Kodaira-vanishing-sea}
\begin{reference}
\cite[21.5.8]{sea}
\end{reference}
Suppose $k$ is a field of characteristic 0,
and $X$ is a
smooth projective $k$-variety. Then for any
ample invertible sheaf $L$,
$H^{i}(X,K_X \otimes L)=0$ for $i>0$.
\end{theorem}

\noindent
For completeness we put another vanishing
theorem
concerning higher-rank vector bundles
(probably due to Serre though Griffith-Harris
don't say so).
The idea is that twisting an ample line
bundle with itself
enough times and tensoring with any vector
bundle $E$
will kill positive cohomology.

\begin{theorem}
\label{theorem-vanis-highe-rank-verct-bund}
\begin{reference}
\cite[p. 159]{gri}
\end{reference}
Let $M$ be a compact, complex manifold and $L
\to M$
a positive line bundle.
Then for any holomorphic vector bundle $E$,
there exists $\mu$ such that
$$
H^q(M,L^\mu \otimes E)=0\qquad \text{for
}q>0, \mu\geq \mu_0.
$$

\end{theorem}

\section{Lefschetz hyperplane sections
theorem}
\label{section-lefschetz-hyperplane}

Lefschetz hyperplane theorem holds for
any $n$-dimensional compact complex manifold
$M$
and $V \subset M$ a smooth hypersurface with
$L=[V]$ positive.
An example of this is when $M$ is projective
and $V$ is the pullback of the hyperplane
bundle
(which, as explained in Section
\ref{section-kodaira-vanishing},
is positive since its curvature form is
Fubini-Study),
that is, $V=M \cap H$ for a generic
hyperplane $H \subset \mathbb{P}^N$.

\begin{theorem}[Lefschetz hyperplane
sections]
\label{theorem-lefschetz-hyperplane}
The map
$$
H^q(M,\mathbb{Q}) \to H^q(V,\mathbb{Q})
$$
induced by the inclusion $V
\xymatrix{\ar@{^{(}->}[r]&} M$
is an isomorphism for $q \leq  n-2$
and injective for $q=n-1$.
\end{theorem}

\begin{proof}
Looks accessible! ``Purely topological''.
\end{proof}

\section{Complex varieties}
\label{section-complex-varieties}

For a definition of sheaf
see Algebraic Geometry Definition
\ref{algebraic-geometry-definition-sheaf}.

\begin{lemma}
\label{lemma-O_X-is-a-sheaf}
Let $X$ be a complex manifold of complex
dimension $n$.

The functor
\begin{align*}
\mathcal{O}_X: \mathit{Open}_X^{\text{op}}
&\longrightarrow \mathit{Set}
\\
U &\longmapsto \mathcal{O}_X(U):=\{f:U\to
\mathbb{C}^n:
f \text{ is holomorphic}\}\\
i:V \hookrightarrow V &\longmapsto
\mathcal{O}_X(i):\mathcal{O}_X(U) \to
\mathcal{O}_X(V)
\end{align*}
where $\mathcal{O}_X(i)$ is given by
restriction,
is a sheaf.
\end{lemma}

\begin{proof}
$\mathcal{O}_X$ is a presheaf,
i.e. a functor,
by definition: restriction
maps the identity to the identity
and preserves compositions.

To check the condition of being a sheaf let
$U_i$ be an open cover of some open set $U$
of $X$
and $f_i:U_i \to X$ sections of
$\mathcal{O}_X(U_i)$,
that is, holomorphic functions, which
coincide pairwise in intersections.
Then it is obvious that the function $f:U \to
\mathbb{C}$,
$x \mapsto f_i(x)$ for any $U_i \ni x$
is well defined. It is also holomorphic
since holomorphicity is a condition defined
on
open sets.
\end{proof}

\medskip\noindent
A meromorphic function is a function that
can be written locally as a quotient of two
holomorphic functions:

\begin{definition}
\label{definition-meromorphic-functions}
\begin{reference}
\cite[Definition 2.1.8]{huc}
\end{reference}
A {\it meromorphic function} on a complex
manifold $X$
is a map
$$
f:X \to \coprod_{x \in
X}\text{Frac}(\mathcal{O}_{X,x})
$$
which associates to any $x\in X$
an element $f_x \in
\text{Frac}(\mathcal{O}_{X,x})$
such that for any $x_0 \in X$
there exists $U\ni x$ open and $f,g
\in\mathcal{O}(U)$
such that $f_x=g/h$ for all $x \in U$.
The {\it sheaf of meromorphic functions} is
$\mathcal{K}_X$.
\end{definition}


\medskip\noindent
Now we pass to the theory of analytic
subvarieties.

\begin{definition}
\label{definition-analytic-subvariety}
An {\it analytic set} or {\it subvariety} $Y$
is a subset of a complex manifold $X$
given locally as the zeros of a finite
collection of holomorphic functions.
\end{definition}

\begin{lemma}
\label{lemma-structure-sheaf-of-subvariety}
Let $Y \subset X$ be an analytic subvariety
of a complex manifold $X$
equipped with the subspace topology.
Then
\begin{align*}
\mathcal{O}_Y: \mathit{Open}_Y^{\text{op}}
&\longrightarrow \mathit{Set} \\
U\cap Y &\longmapsto \{f|_{U \cap Y}:f \in
\mathcal{O}_X(U)\}
\end{align*}

\noindent
is a sheaf.
\end{lemma}

\begin{proof}
This is a presheaf for the same reasons as
the sheaf $\mathcal{O}_X$ is a presheaf.
The fact that it is a sheaf follows
from the fact that $\mathcal{O}_X$ is a
sheaf and by definition of subspace topology;
namely for any family of functions
defined over an open cover $\{U_i \cap Y\}$
of some open set  $U\cap Y$ we can
reconstruct a function
on $U$ that restricted to $U \cap Y$ is as
desired.
\end{proof}

Recall that any map of sheaves of groups
has well-defined kernel, see Lemma
\ref{algebraic-geometry-lemma-sheaves-%
valued-on-groups-have-kernels}.

\begin{definition}
\label{definition-ideal-sheaf}
Let $Y \subset X$ be an analytic subvariety
of a complex manifold $X$.
The {\it ideal sheaf} of $Y$, denoted
$\mathcal{I}_Y$,
is the kernel
of the restriction map $\mathcal{O}_X \to
\mathcal{O}_Y$.
\end{definition}

\noindent
This just says that $\mathcal{I}_X(U)$ is the
set of functions $f \in \mathcal{O}_X$ that
vanish under the restriction map, i.e.
that vanish at points of $Y$.

Since the restriction map is surjective,
we have the following short exact sequence:
$$
\xymatrix{
0\ar[r]&\mathcal{I}_Y\ar[r]&\mathcal{O}
_X\ar[r]&\mathcal{O}_Y\ar[r]&0
}
$$

\medskip\noindent
Notice that for any open set $U\subset X$,
$\mathcal{I}_Y(U)$ is an ideal of
$\mathcal{O}_X(U)$. (Indeed,
if a function vanishes at $Y$ then its
product
with any other function will also vanish.)

\begin{definition}
\label{definition-irreducible-variety}
An analytic variety $V\subset U\subset
\mathbb{C}^n$ is {\it irreducible} if
$V$ cannot be written as the union of two
distinct analytic varieties
$V_1,V_2\subset U$, both distinct to $V$.
\end{definition}

\medskip\noindent
We will show that for the case of an
irreducible hypersurface
the ideal sheaf is locally generated by a
single element.
This will allow us to construct a line bundle
from the defining functions
of the hypersurface, which is in general not
possible
for an arbitrary smooth hypersurface,
--- although it suffices to ask for
orientability in the
smooth case.

We need two ingredients:
(1) the fact that the stalks $\mathcal{O}_p$
of the
sheaf of holomorphic functions are UFDs,
and (2) the following lemma.

\begin{lemma}
\label{lemma-ideal-irred-hyper-heigh-1}
The ideal
(correction: this should be the stalk)
of an irreducible analytic hypersurface is a
height-1 prime ideal.
\end{lemma}

\begin{proof}
We already know that irreducible varieties
correspond with prime ideals.

To show the ideal $I$ defining the variety
$Y$ has height 1
suppose that $0 \subset \mathfrak{p}
\subseteq I$ is prime.
Then $V(\mathfrak{p})$ is an analytic variety
that contains $Y$.
At regular points of both $V(\mathfrak{p})$
and $Y$,
both are smooth manifolds, but since $Y$ is
of codimension 1 and
$V(\mathfrak{p})$ does not equal all of  $Y$,
we conclude that they coincide.
Thus, in a neighbourhood
of a regular point we have
$\mathfrak{p}=I(V(\mathfrak{p}))=I$ by the
Nullstellensatz.
\end{proof}

\begin{lemma}
\label{lemma-stalk-ideal-sheaf-princ-hyper}
The stalks of the ideal sheaf of
an irreducible hypersurface
are principal rings.
\end{lemma}

\begin{proof}
Since $\mathcal{O}_p$
is a UFD, we pick any element
$f \in \mathcal{I}_p$, and factor it as a
product of irreducibles.
Since $\mathcal{I}$ is prime, some of the
irreducible factors,
say $f_i$, is in $\mathcal{I}_p$.
The ideal generated by $f_i$ is prime since
$\mathcal{O}_p$
js a UFD (see
\href{https://stacks.math.columbia.edu/tag/%
034T}{034T}).
But by Lemma
\ref{lemma-ideal-irred-hyper-heigh-1},
$\mathcal{I}_p$ is a height-1 prime ideal, so
the prime ideal $(f_i)$
must equal all of $\mathcal{I}_p$.
\end{proof}

\begin{lemma}
\label{lemma-ideal-sheaf-irred-hyper-line}
The ideal sheaf $\mathcal{I}$
of the irreducible hypersurface $Y \subset X$
corresponds to a unique line bundle.
\end{lemma}

\begin{proof}
Let $p \in U$ and consider for any stalk
$\mathcal{I}_p$
its generator $[f_p]$,
which has a representative $(f_p,U_p)$.
These form an open cover of $X$,
coincide up to a unit in the intersections as
follows.
These units shall give us a cocycle.

For any $U_i\cap U_j \neq 0$
we have
[This is slightly imprecise.
We may extend $f_j|_{U_i\cap U_j}$ to $U_i$
and $f_i|_{U_i\cap U_j}$ to $U_j$,
and then restrict back to $U_i\cap U_j$.
This works because we can extend holomorphic
functions by their local power series
expansion!],
$f_j\in (f_i)$, so there exist functions
$g_{ij}$ such that
$$
f_i=g_{ij}f_j.
$$
These satisfy the cocycle condition since
$$
f_if_jf_k=g_{ij}f_jg_{jk}f_kg_{ki}f_i.
$$
By Proposition
\ref{proposition-vecto-bundl-cocyc} we have a
line
bundle.
\end{proof}

\noindent
Finally we have concluded that the ideal
sheaf
of a hypersurface is a line bundle. Now we
can take its dual to define:

\begin{definition}
\label{definition-sheaf-assoc-effec-carti}
Let $D:=S$ be an effective Cartier divisor,
that is,
an analytic hypersurface of a complex
manifold.
The line bundle $\mathcal{O}(D)$ is the dual
line bundle
of the ideal sheaf of $D$.
\end{definition}

\section{Chow's theorem}
\label{section-chow-theorem}

\begin{theorem}[Chow]
\label{theorem-chow}
Any analytic subvariety of $\mathbb{C}P^{n}$
is algebraic.
\end{theorem}

\begin{proof}
First we prove it for hypersurfaces.
A hypersurface in $\mathbb{P}^n$ defines a
line bundle
along with a section vanishing on the
hypersurface.
Such a line bundle must be $\mathcal{O}(d)$
for some $d$
by Exercise \ref{exercise-pic-of-pn-is-z}.
Sections of $\mathcal{O}(d)$ are homogeneous
degree-$d$ polynomials
by Theorem
\ref{theorem-secti-od-degre-d-homog-poly}.

For higher codimension the challenge is find
adequate planes
in $\mathbb{P}^n$, so that we can project
$V$, $\dim V=k$
to $\mathbb{P}^{k+1}$ and apply the
hypersurface version.
Pick any point $p\in \mathbb{P}^n$ not in
$V$.
We can choose a $(n-k-1)$-plane not
intersecting $V$ and containing $p$.
Inside such $\mathbb{P}^{n-k-1}$ we can pick
a $\mathbb{P}^{n-k-2}$
not touching $p$.
Now pick a complementary $\mathbb{P}^{k+1}$
in $\mathbb{P}^n$,
i.e. $\mathbb{P}^n=\mathbb{P}^{k+1}\oplus
\mathbb{P}^{n-k-2}$.

Now project $V$ to $\mathbb{P}^{k+1}$ and
apply Chow for hypersurfaces.
We obtain a polynomial $F$ defining the
projection of $V$.
Extend $F$ to $\tilde{F}$ on all
$\mathbb{P}^n$ by the rule
$\tilde{F}(x_0,…,x_n)=F(x_0,…,x_{k+1})$.
Note that since $p$ is not in $V$,
$\tilde{F}$ does not vanish in $p$.
Thus $V$ is the intersection of all such
polynomials $F$,
one for each $p$, which ensures that no
points other than those
of $V$ are in the common zeroes.
\end{proof}

\section{Kodaira Embedding theorem}
\label{section-kodaira-embedding}

\noindent
We introduced line systems in Section
\ref{section-line-bundles}.
Let $L$ be a holomorphic bundle on a complex
manifold $X$.
A point $x \in X$ is
a {\it base point} of $L$ if $s(x)=0$ for all
$s\in H^{0}(X,L)$
and the {\it base locus} $\text{Bs}(L)$
is the set of all base points of $L$.

Let $s_0,\ldots,s_N\in H^{0}(X,L)$ are a
basis
for some subspace $E \subset H^0(X,L)$.
Define a map
\begin{align*}
\varphi_L: X\setminus\text{Bs}(L)
&\longrightarrow \mathbb{P}^N \\
x &\longmapsto [s_0(x):\ldots:s_N(x)]
\end{align*}

\noindent
which is well-defined since we excluded
the base locus,
and holomorphic since we chose holomorphic
sections of $L$ for the coordinates.
Further, the pullback of the hyperplane
bundle of $\mathbb{P}^N$ is $L$,
i.e.
$$
\varphi^*_L \mathcal{O}_{\mathbb{P}^N}(1)
\simeq L|_{X\setminus\text{Bs}(L)}
$$
since $s_i=0$ is the equation of a hyperplane
in $\mathbb{P}^N$,
so, it is a section of $H$.
Its pullback under $\varphi$ gives the
section $s_i$
whose associated divisor $(s_i)$
corresponds to the line bundle $L$ ---
the line bundle associated to a divisor
associated to a section of the line bundle
is the line bundle itself.

\medskip\noindent
The question is: when is $\varphi$ an
embedding?

\begin{definition}
\label{definition-ample}
\begin{reference}
\cite[Definition 2.3.28]{huc}
\end{reference}
A line bundle $L$ on a complex manifold is
called {\it ample} if for some $k>0$
and some linear system in $H^{0}(X,L^k)$ the
associated map $\varphi$ is an
embedding.
\end{definition}

\begin{theorem}[Kodaira Embedding]
\label{theorem-Kodaira-embedding}
\begin{reference}
\cite[Theorem 10.12]{lec}
\end{reference}
Suppose $M$ is a compact complex manifold.
A holomorphic line bundle $L\to M$ is
ample if and only if it is positive.
Thus $M$ is projective if and only if it
admits a positive holomorphic line bundle.
\end{theorem}

\begin{proof}
We must show that $\varphi$ separates points
and has
surjective differential everywhere…
\end{proof}

\noindent
Compare this theorem with Nakai-Moishezon
Criterion
\ref{algebraic-geometry-theorem-%
Nakai-Moishezon-criterion}.

\begin{exercise}
\label{exercise-l-ample-impli-lotim-globa}
Let $L$ be an ample bundle on a K3 surface
$M$.
Prove that $L^{\otimes 2}$ is
globally generated
(that is, for each $x \in M$ there exists a
section $h \in
H^{0}(L^{\otimes 2})$ which does not vanish
in $x$).
\end{exercise}

\begin{proof}
\begin{reference}
\href{https://math.stackexchange.com/%
questions/5081481/
ample-line-bundle-on-a-k3-surface-implies-l-
otimes-2-is-base-point-free/
5081492?noredirect=1#comment10928708_5081492}
{StackExchange}
\end{reference}
As an outline for exposition:
\begin{enumerate}
\item Show that $L$ has sections
(by ampleness+Kodaira Vanishing+Riemann-Roch
for surfaces),
pick a section $C$ and note that if
$L^{\otimes2}$ had base points they would be
in $C$.

\item  Using Kodaira Vanishing again we show
that every section of the line bundle
restricted to $C$
comes from a section of the line bundle on
$M$.
Thus it's enough to find a nonvanishing
section on the restricted bundle.

\item Then we aim to apply Riemann-Roch on
$C$.
If we can show that $h^0((L^{\otimes 2})^\vee
\otimes \omega_C)$
and $h^0((L^{\otimes 2}(-p))^\vee \otimes
\omega_C)$
both vanish, we apply Riemann-Roch and see
that
$h^0(L^{\otimes 2})> h^0(L^{\otimes 2}(-p)$,
which proves the result since that says
that not every section is a section vanishing
at $p \in C$.

\item To compute the vanishings of the $h^0$
in the previous step we show they have
negative degrees
(genus formula + degree of canonical bundle
of a curve).
\end{enumerate}

First we need to show that $L$ has sections.
To see this I will show that since $L$ is
ample, $L^2>0$. Indeed, if $mL$ is very ample
we apply by Riemann-Roch, Kodaira Vanishing
and the fact that $\chi(\mathcal{O}_X)=2$ on
a K3 to obtain $h^0(mL)=2+\frac{1}{2}m^2L^2$
which must be positive since $mL$ has
sections by being very ample. Notice $L^2$
cannot be zero as $h^0(mL)=2$ would imply we
have an embedding of a K3 surface into
$\mathbb{P}^1$.

Applying again Riemann-Roch and Kodaira
vanishing we see that $h^0(L)>0$. Let $C$ be
a section of $L$.

Notice that if $L^{\otimes 2}$ had any base
points, they would have to be on $C$. Indeed,
since $C$ is the vanishing set of a section
$s\in H^{0}(L)$, we must have $s \otimes s
\in H^{0}(L^{\otimes 2})$ vanishing in any
base point of $L^{\otimes 2}$.

Then we consider the ideal exact sequence for
$C$ and tensor by $L^{\otimes 2}$ (which
preserves left-exactness because line bundles
are locally free of rank 1, that is, we are
tensoring by the underlying ring) to obtain
$$ \xymatrix{ 0\ar[r]&L^{\otimes
2}(-C)\ar[r]&L^{\otimes 2}\ar[r]& L^{\otimes
2}\otimes\mathcal{O}_C\ar[r]&0 } $$ Notice
that the term on the left is actually $L$
since the ideal sheaf
$\mathcal{O}_M(-C)=\mathcal{I}_C$ is dual to
$L$ because $C$ is defined as the vanishing
set of a section of $L$. Passing to
cohomology we obtain $$ \xymatrix{
H^{0}(L^{\otimes 2})\ar[r]&H^{0}(L^{\otimes
2}\otimes \mathcal{O}_C)\ar[r]& H^{1}(L) } $$
where the latter term vanishes by Kodaira
Vanishing because $L$ is ample. This says
that every section of $L^{\otimes 2}|_C$ is
the restriction to $C$ of a section of
$L^{\otimes 2}$. In turn this shows that it's
enough to show that the bundle $L^{\otimes
2}$ has no base points along $C$.

Now pick a point $p\in C$. Saying
that $p$ is not a base point of $L^{\otimes
2}|_{C}$ is the same as saying that not every
section of $L^{\otimes 2}|_{C}$ vanishes at
$p$, that is, that $h^0(L^{\otimes
2}|_{C}(-p))<h^0(L^{\otimes 2}|_{C})$,
which concludes the exercise.

The key step to prove
the last inequality is showing that
$h^0(\omega_C\otimes(L^{\otimes2}|_{C})^\vee)
=0=h^0(\omega_C\otimes
(L^{\otimes2}|_{C}(-p))^\vee)$. For this we
need to see these line bundles as Weil
divisors, so that their degree is the sum of
their multiplicities. First note that by the
genus formula for (possibly singular) curves
on smooth surfaces and by the fact that $M$
is smooth, we see that $2p_a(C)-2=(L.L)$ so
that
$$
4p_a(C)-4=(2L.L)=\text{deg}_C(L^{\otimes 2})
$$
Notice that since $L^2>0$ we exclude the
cases that $p_a(C)=0,1$.

By Riemann-Roch on the curve
$C$, we know that
$\text{deg}\omega_C=2p_a-2$. Then
$$
\text{deg}_C
(\omega_C\otimes (L^{\otimes2})^\vee)
=\text{deg}_C(\omega_C)-\text{deg}
_C(L^{\otimes
2})
=2p_a-2-(4p_a-4)<0
$$
and a similar computation works for
$L^{\otimes 2}(-p)$ as it has the same
degree of $L^{\otimes 2}$ minus 1. This
implies that neither of these bundles
can have sections since any section would
provide a linearly equivalent
effective divisor
--- degree is preserved under linear
equivalence,
and divisors given by sections are effective.
Therefore
$h^0(\omega_C\otimes
(L^{\otimes2}|_{C})^\vee)=0
=h^0(\omega_C\otimes (L^{\otimes
2}|_{C}(-p))^\vee)$ as claimed.

Finally we apply Riemann-Roch to $L^{\otimes
2}$ and $L^{\otimes 2}(-p)$ to
obtain that $h^0(L^{\otimes
2})>h^0(L^{\otimes 2}(-p))$.
\end{proof}

\section{Riemann-Roch formulas}
\label{section-riemann-roch}

\noindent
We want to know the dimension of
$H^0(X,\mathcal{O}(D))$,
where $X$ is a curve and $D$ is a divisor on
$X$.
To refresh out memory recall that if $D$ is a
smooth hypersurface
(in this case this just means that it's a set
of points
without multiplicities)
then that $\mathcal{O}(D)$ is given by the
cocycle
$g_{ij}=f_j/f_i$ where $D=\text{div}(f_i)$ on
$U_i$.
The $g_{ij}$ give a cocycle, which in turn is
a line bundle on $X$.
If $e_i$ is a local basis then we have
$e_i=g_{ij}e_j$
by definition of transition function.
A section is given locally by $h_ie_i$ and it
has to glue,
that is, $h_je_j=h_ie_i=e_if_j/f_ie_j$
so $h_j/f_j=h_i/f_i$.
That's literally the data of a meromorphic
function,
namely some local meromorphic functions which
coincide in intersections.

What's the vanishing of this meromorphic
function along $D$?
Now suppose $D$ is an arbitrary Weil divisor,
so the $f_i$ have some multiplicities.
If $h_i$ has no zeroes, then the multiplicity
of the new
meromorphic function remains unaltered, but
$h_i$ can also vanish
along $D$ and then the vanishing becomes as
positive as we want it.
But it cannot have poles, the least positive
it can be is
the prescribed negativity of $D$.

This is why looking for $h^0(D)$ is just as
much as looking for
meromorphic functions $h$ on $X$ such that
$$
\text{div}h + D \geq 0\qquad \text{(is
effective).}
$$



\begin{theorem}[Riemann-Roch for curves]
\label{theorem-Riemann-Roch-for-curves}
Let $D$ be a divisor on a compact Riemann
surface $X$, that is, $D$ is a
collection of points with multiplicities on
$X$.
\begin{equation}
\label{equation-riema-roch-riema-surfa}
h^0(D)-h^0(K-D)=d-g+1
\end{equation}
\end{theorem}

\begin{exercise}
\label{exercise-no-holom-vecto-field-g}
Prove that if $X$ is a complex smooth curve
with $g\geq 2$ then $h^0(T_X)=0$.
\end{exercise}

\begin{proof}
First you notice that the dimension of the
holomorphic 1-forms is the genus.
This is just because the genus $p_a$ is
defined as $h^1(\mathcal{O}_X)$, which
is just $h^0(K_X)$ by Serre duality. Then you
say well if you have a holomorphic
vector field, pair it with the nonzero
holomorphic 1-form. This gives a nonzero
function, but it must be constant because $X$
is compact. Then it is actually
nowhere vanishing. This says the vector field
cannot vanish anywhere, which
means the tangent bundle of the curve is
trivial. Apparently elliptic curves,
i.e. $g=1$ are the
\end{proof}

\begin{theorem}[Riemann-Roch for line bundles
on surfaces]
\label{theorem-riema-roch-line-bundl-surfa}
Let $L$ be a line bundle on a complex surface
$X$, and $K_X=\Omega^2X$ the
canonical bundle of $X$. Then
\begin{equation}
\label{equation-riema-roch-line-bundl-surf}
\chi(L)=\chi(\mathcal{O}_X)+\frac{1}{2}(L,L-
K_X)
\end{equation}
\end{theorem}

\begin{lemma}[Genus formula]
\label{lemma-genus-formula}
Let $C$ be any curve on a surface $S$. Then
\begin{equation}
\label{equation-genus-formula}
2p_a(C)-2=(C.C+K_S).
\end{equation}
\end{lemma}

\begin{proof}
By Remark
\ref{remark-euler-characteristic-additive}
below we can take Euler characteristic
of the ideal sheaf exact sequence of $C$
to obtain
$$
\chi(\mathcal{O}(-C))+\chi(\mathcal{O}_C)=
\chi(\mathcal{O}_S).
$$
Since $p_a(C)=h^0(K_S)
=1-\chi(\mathcal{O}_C)$, we
obtain that
$p_a(C)=1-\chi(\mathcal{O}_S)-
\chi(\mathcal{O}(-C))$.
To compute $\chi(\mathcal{O}(-C))$ we use
Riemann-Roch formula for curves on surfaces
to get
$$
\chi(\mathcal{O}(-C))=\chi(\mathcal{O}_S)
+\frac{(\mathcal{O}(-C).
\mathcal{O}(-C)-K_S)}{2}
$$
Taking duals on both entries of the
intersection form and substituting in our
previous expression for $p_a(C)$ we obtain
the result.
\end{proof}

\begin{remark}
\label{remark-euler-characteristic-additive}
To prove that Euler characteristic is
additive, i.e. that for an exact sequence
$$
\xymatrix{
0\ar[r]
&E\ar[r]
&F\ar[r]
&G\ar[r]
&0
}
$$
we must have
$$
\chi(F)=\chi(E)+\chi(G),
$$
we take the long exact sequence in cohomology
and use the general fact that for
a finite exact sequence of vector spaces
$$
\xymatrix{
0\ar[r]
&V_0\ar[r]
&V_1\ar[r]
&\cdots\ar[r]
&V_n\ar[r]
&0
}
$$
we must have
$$
\sum_i(-1)^i\dim V_i,
$$
which basically follows by dimension theorem.
\end{remark}

\begin{exercise}
\label{exercise-singular-curve-on-k3}
Let $C$ be a singular, irreducible complex
curve on a K3 surface $X$. Prove that
$(C.C)\geq 0$.
\end{exercise}

\begin{proof}
This is just an application of Eq.
\ref{equation-genus-formula}. Since
$K_X=\mathcal{O}_X$ it suffices to show that
the arithmetic genus
$p_a(C)$ is not zero. This follows from the
fact that $C$ is singular, since its
arithmetic genus, defined as 
$p_a(C):=1-\chi(\mathcal{O}_C)$,
satisfies
$$
p_a(C)=g(C) + h^0(\delta)
$$
where $g(C)$ is the geometric genus of $C$,
defined as
the genus of the normalization $\tilde{C}$,
i.e. $g(C):=g(\tilde{C})$,
while $\delta:=\nu_*\mathcal{O}_{\tilde{C}}
/\mathcal{O}_C$
where $\nu:\tilde{C} \to C$ is the
normalization map.

So we have to prove that equation
and also prove that $\delta$
always has sections…
\end{proof}


\section{Symplectic isogenies and
Shafarevich}
\label{section-symplectic-isogenies-shaf}

These are notes from Misha Verbitsky's talk
{\it Symplectic isogenies for hyperkahler
manifolds and Shafarevich conjecture},
given at IMPA on May 7, 2026.

\medskip
Abstract.
A correspondence between hyperkahler
manifolds is a symplectic isogeny if it is
generically finite over both factors and its
push-pull operation sends the holomorphic
symplectic form to a nonzero holomorphic
symplectic form.
The conjecture, attributed to Shafarevich by
Buskin, says that hyperkahler manifolds in
the same deformation family with isomorphic
rational period lattices should be related
by such an isogeny.
The talk explains degenerate twistor
families and shows that projective
hyperkahler manifolds in the same such
family are symplectically isogenous.

\medskip
Absolute upshot.
Same degenerate twistor family is a much
more structured condition than being merely
deformation-equivalent.
Misha proves the Shafarevich-type isogeny
statement first in this structured setting:
the common Lagrangian fibration gives torus
fibers and torsors, which can be trivialized
after finite base change.
The resulting graph produces the desired
symplectic correspondence.

\medskip
Upshot.
Find the simplest Hodge-theoretic piece
forced by the holomorphic symplectic form:
the transcendental lattice $V^{tr}$.
Then prove that this piece behaves well:
it is usually irreducible, and in the
degenerate case has only one controlled
defect.
See Theorem
\ref{theorem-transcendental-k3-type}.

\medskip
{\bf Motivation: ADE contractions.}

Let $X$ be a smooth K3 surface and let
$C \subset X$ be a smooth rational curve.
Since $K_X=\mathcal{O}_X$, adjunction gives
$$
2g(C)-2=C^2+K_X\cdot C=C^2.
$$
Thus a smooth rational curve on a K3 surface
has self-intersection $C^2=-2$.
Conversely, a smooth curve with $C^2=-2$
has genus $0$.
Such a curve is called a $(-2)$-curve.

\medskip
A $(-2)$-curve is not obtained by blowing up
another smooth K3 surface.
Rather, it can be contracted:
$$
X \longrightarrow Y.
$$
The target $Y$ is usually singular.
For a single $(-2)$-curve, the image point is
a Du Val singularity of type $A_1$.
Locally this is the rational double point
$$
xy=z^2,
$$
or equivalently $\mathbb{C}^2/\{\pm 1\}$.

\medskip
More generally, suppose
$C_1,\ldots,C_r \subset X$ are $(-2)$-curves
whose intersection matrix is negative
definite.
Then the whole configuration can be
contracted.
The dual graph has one vertex for each
$C_i$ and an edge when two curves meet
transversely once.
For rational double points this graph is one
of the ADE Dynkin diagrams:
$$
A_n,\qquad D_n,\qquad E_6,E_7,E_8.
$$

\medskip
Conversely, the minimal resolution of a
Du Val singularity replaces the singular
point by a configuration of $(-2)$-curves.
Their intersection graph is the corresponding
Dynkin diagram.
Thus the slogan is
$$
\text{Du Val singularity}
\longleftrightarrow
\text{ADE configuration of }
(-2)\text{-curves}.
$$
The lattice generated by these curves is the
corresponding root lattice.

\medskip
In the motivating picture, one contracts a
Dynkin configuration
$$
f:X\longrightarrow Y
$$
and then resolves the singular surface $Y$:
$$
\widetilde{Y}\longrightarrow Y.
$$
If $X$ is not already the chosen resolution,
the common singular surface $Y$ gives a
relation between $X$ and $\widetilde{Y}$.
One studies this relation by the
correspondence
$$
X \xleftarrow{\ p_1\ }
X\times_Y\widetilde{Y}
\xrightarrow{\ p_2\ }\widetilde{Y}.
$$
This is the prototype for replacing a
birational map by a cycle in a product.

\medskip
This is the end of the motivating example.
The real subject of the talk is to replace
maps between hyperkahler manifolds by
correspondences which transport holomorphic
symplectic forms.

\medskip
{\bf Correspondences and symplectic
isogenies.}

\begin{definition}
\label{definition-correspondence-verbitsky}
Let $M$ be a complex variety.
A {\it cycle} in $M$ is a subvariety
together with a positive integer
multiplicity on each irreducible component.
Let $A,B$ be irreducible compact complex
varieties of dimension $n$.
A {\it correspondence} from $A$ to $B$ is an
$n$-dimensional cycle
$$
Z\subset A\times B
$$
such that every irreducible component of
$Z$ projects surjectively to both $A$ and
$B$.
\end{definition}

\noindent
Thus a correspondence is a multivalued map.
When using cohomology, we often think of the
cycle $Z$ as its Poincare dual differential
form. The caveat is that pushforward of a
differential form is generally a current.
For holomorphic forms in this setting, the
pushforward is holomorphic away from the
exceptional locus and extends by Hartogs.

\begin{definition}
\label{definition-symplectic-isogeny}
Let $A,B$ be holomorphically symplectic
manifolds and let
$Z\subset A\times B$ be a correspondence.
Let $\pi_1,\pi_2$ be the projections from
$Z$ to $A$ and $B$.
The correspondence $Z$ is a
{\it symplectic isogeny} if, for a
holomorphic symplectic form $\Omega_A$ on
$A$,
$$
(\pi_2)_*(\pi_1^*\Omega_A)\neq 0.
$$
\end{definition}

\begin{lemma}
\label{lemma-sympl-isog-transcendental}
Let $M_1,M_2$ be holomorphically symplectic
manifolds of K3 type, and let
$Z\subset M_1\times M_2$ be a symplectic
isogeny.
Then the push-pull map
$$
\Phi_Z(\eta):=(\pi_2)_*(\pi_1^*\eta)
$$
induces an isomorphism of rational
transcendental Hodge lattices
$$
\Phi_Z:T(M_1)\longrightarrow T(M_2).
$$
\end{lemma}

\noindent
Here $T(M)$ means the transcendental Hodge
lattice $V^{tr}\subset H^2(M,\mathbb{Q})$.
This is not literally the complex torus
attached to a weight-one Hodge structure.
The torus story is the model case: isogenies
of tori are the same as isomorphisms of
rational weight-one Hodge structures.
For K3-type Hodge structures, the analogous
object is $T(M)$.

\medskip
Recall also the dictionary between complex
tori and Hodge structures; see
{\it mumford-tate-groups-in-hodge-theory.tex}
for the Hodge-theoretic background.
If $T$ is a compact complex torus, then
$H_1(T,\mathbb{Z})$ carries an integral
Hodge structure of weight $1$.
Conversely, such a Hodge structure
$$
V_{\mathbb{C}}=V^{1,0}\oplus V^{0,1}
$$
with a lattice $V_{\mathbb{Z}}\subset V$,
gives the complex torus
$$
T=V^{1,0}/V_{\mathbb{Z}}.
$$
After tensoring the lattice with
$\mathbb{Q}$, this is the category of
complex tori up to isogeny.

\begin{definition}
\label{definition-isogeny-tori}
Let $T_1,T_2$ be compact complex tori of the
same dimension.
An {\it isogeny} from $T_1$ to $T_2$ is a
subtorus
$$
T\subset T_1\times T_2
$$
such that the two projections
$T\to T_1$ and $T\to T_2$ are coverings.
We say that $T_1$ and $T_2$ are isogenous
if such a subtorus exists.
\end{definition}

\medskip
Briefly, the Mumford-Tate group records the
rational algebraic symmetries forced by a
Hodge structure.
See Section
\ref{mumford-tate-section-mumford-tate-group}
and Definition
\ref{mumford-tate-definition-mumford-tate-%
group}.
If $V$ is a rational Hodge structure and
$$
\rho:U(1)\longrightarrow
GL(V_{\mathbb{R}})
$$
is the Hodge circle action, then $MT(V)$ is
the smallest algebraic subgroup of $GL(V)$
defined over $\mathbb{Q}$ whose real points
contain $\rho(U(1))$.

\begin{proposition}
\label{proposition-mt-generated-galois}
Let $V$ be a rational Hodge structure.
For each
$\sigma\in\operatorname{Aut}_{\mathbb{Q}}
(\mathbb{C})$,
let $U(1)^\sigma$ denote the Galois conjugate
of the Hodge circle in
$GL(V_{\mathbb{C}})$.
Then $MT(V)_{\mathbb{C}}$ is the Zariski
closure of the subgroup generated by all
$U(1)^\sigma$.
\end{proposition}

\begin{theorem}
\label{theorem-transcendental-hodge-lattice}
Let $V_{\mathbb{Q}}$ be a rational Hodge
structure of weight $d$.
Let $V^{tr}\subset V_{\mathbb{Q}}$ be the
smallest rational subspace such that
$$
V^{d,0}\subset V^{tr}_{\mathbb{C}}.
$$
Equivalently, $V^{tr}_{\mathbb{C}}$ is
generated by all Galois conjugates of
$V^{d,0}$.
Then $V^{tr}$ is a Hodge substructure.
\end{theorem}

\begin{definition}
\label{definition-polarization-hodge}
A {\it polarization} on a rational Hodge
structure $V$ of weight $k$ is a
$U(1)$-invariant nondegenerate bilinear form
$$
h\in V_{\mathbb{Q}}^*\otimes
V_{\mathbb{Q}}^*
$$
which is symmetric or skew-symmetric
depending on the parity of $k$, and which
satisfies the Hodge-Riemann positivity
condition.
\end{definition}

\noindent
This is Misha's formulation.
In
Definition
\ref{mumford-tate-definition-polarization},
Andrey packages the same positivity using a
morphism to $\mathbb{Q}(-k)$ and the Weil
operator.

\medskip
The basic source of such polarizations is the
Hodge-Riemann pairing.
For a compact Kahler manifold $X$, a Kahler
class $\omega$ gives, on primitive
$H^k(X)$, the pairing
$$
(\alpha,\beta)\longmapsto
\int_X\alpha\wedge\beta\wedge\omega^{n-k}.
$$
After the usual sign and Weil-operator
normalization, the Hodge-Riemann bilinear
relations say this pairing is positive.
Thus the abstract polarization in Hodge
theory is the linear-algebra shadow of the
intersection pairing from complex geometry.

\begin{definition}
\label{definition-simple-semisimple-object}
Let $\mathcal{A}$ be an abelian category.
An object $V$ is {\it simple} if its only
subobjects are $0$ and $V$.
It is {\it semisimple} if it is a finite
direct sum of simple objects.
\end{definition}

\noindent
The upshot is that polarizations prevent
extensions from hiding inside a Hodge
structure.
If $W\subset V$ is a Hodge substructure, the
orthogonal complement $W^\perp$ is again a
Hodge substructure, so $V$ splits as
$W\oplus W^\perp$.

\begin{proposition}
\label{proposition-polarizable-hs-semisimple}
Every polarizable rational Hodge structure
is semisimple.
\end{proposition}

\noindent
In particular, the rational Hodge structures
$H^k(X,\mathbb{Q})$ of a compact Kahler
manifold are semisimple.
Indeed, the Lefschetz decomposition reduces
the statement to primitive cohomology, where
the Hodge-Riemann bilinear relations give
the required polarization.

\begin{definition}
\label{definition-k3-type-hodge-structure}
A rational Hodge structure of weight $2$
is of {\it K3 type} if
$$
V_{\mathbb{C}}=
V^{2,0}\oplus V^{1,1}\oplus V^{0,2}
$$
and $\dim V^{2,0}=1$.
\end{definition}

\begin{example}
\label{example-bbf-k3-type}
Let $M$ be an irreducible holomorphic
symplectic manifold.
Then $H^2(M,\mathbb{Q})$ is a Hodge
structure of K3 type.
It carries the Beauville-Bogomolov-Fujiki
quadratic form
$$
q:H^2(M,\mathbb{Q})\times H^2(M,\mathbb{Q})
\longrightarrow \mathbb{Q}.
$$
Fujiki's formula says that, if
$\dim_{\mathbb{C}}M=2n$, then
$$
\int_M\alpha^{2n}=c\,q(\alpha,\alpha)^n
$$
for a positive rational constant $c$.
\end{example}

\noindent
Thus, for hyperkahler manifolds, the BBF
form plays the role of the natural quadratic
form on the K3-type Hodge structure
$H^2(M)$.
It is not literally the ordinary
Hodge-Riemann pairing on all cohomology, but
it is the replacement on $H^2$ which controls
the period lattice and the Torelli theory.

\begin{theorem}
\label{theorem-transcendental-k3-type}
Let $(V,q)$ be a pseudo-polarized rational
Hodge structure of K3 type.
Let $V^{tr}\subset V$ be the transcendental
Hodge sublattice, i.e. the smallest rational
Hodge substructure containing $V^{2,0}$.
Then exactly one of the following happens.
\begin{enumerate}
\item If $V^{1,1}_{\mathbb{Q}}$ contains a
vector $x$ with $q(x,x)>0$, then
$q|_{V^{tr}}$ is a polarization and
$V^{tr}$ is irreducible.
\item If $q|_{V^{1,1}_{\mathbb{Q}}}$ is
negative definite, then $V^{tr}$ is
irreducible and
$q|_{V^{tr}}$ is a pseudo-polarization.
\item If $q|_{V^{1,1}_{\mathbb{Q}}}$ is
negative semidefinite with one-dimensional
kernel, then
$$
V^{1,1}_{\mathbb{Q}}\cap
(V^{1,1}_{\mathbb{Q}})^\perp
=\langle\eta\rangle,
\qquad q(\eta,\eta)=0.
$$
In this case $\langle\eta\rangle\subset
V^{tr}$ is the only nonzero proper Hodge
substructure of $V^{tr}$, and
$V^{tr}/\langle\eta\rangle$ is irreducible.
\end{enumerate}
\end{theorem}

\noindent
The third case is the strange one.
The transcendental lattice is no longer
irreducible, but it fails in the smallest
possible way: it has a unique isotropic
rational line, and the quotient by this line
is irreducible.

\medskip
Upshot: for K3-type Hodge structures the
transcendental part is almost simple.
Usually it is irreducible.
In the degenerate case it has exactly one
visible defect, the isotropic line
$\langle\eta\rangle$.
Thus a correspondence which is nonzero on
$V^{2,0}$ is forced to control essentially
all of $V^{tr}$.
This is why the condition
$(\pi_2)_*\pi_1^*\Omega\neq 0$ is so strong.

\begin{lemma}
\label{lemma-sympl-isog-controls-tm}
The point of the symplectic isogeny
condition is that the correspondence is
nonzero on the line $H^{2,0}$.
Since $T(M)$ is the smallest rational Hodge
substructure generated by this line, the
correspondence is forced to identify the
transcendental lattices:
$$
T(M_1)\cong T(M_2).
$$
\end{lemma}

\begin{theorem}[Witt]
\label{theorem-witt-quadratic-spaces}
Let $(V,q)$ be a nondegenerate rational
quadratic space.
Let $W_1,W_2\subset V$ be subspaces such that
$q|_{W_i}$ is nondegenerate.
If $W_1$ and $W_2$ are isomorphic as
quadratic spaces, then
$W_1^\perp$ and $W_2^\perp$ are isomorphic
as quadratic spaces.
\end{theorem}

\begin{lemma}
\label{lemma-isogeny-gives-h2-hodge-isom}
Let $M_1,M_2$ be projective hyperkahler
manifolds in the same deformation class.
If $M_1$ and $M_2$ are symplectically
isogenous, then
$$
H^2(M_1,\mathbb{Q})\cong
H^2(M_2,\mathbb{Q})
$$
as rational Hodge structures.
\end{lemma}

\noindent
Indeed, the isogeny identifies the
transcendental Hodge lattices $T(M_i)$.
The algebraic parts are the orthogonal
complements with respect to the BBF form.
Witt's theorem then gives the corresponding
isomorphism on the complements, hence on all
of $H^2$.

\begin{theorem}[Torelli for K3 surfaces]
\label{theorem-torelli-k3}
Let $S_1,S_2$ be K3 surfaces.
Then $S_1\cong S_2$ if and only if their
integral Hodge structures on $H^2$ are
isomorphic, with the Kahler cone condition.
\end{theorem}

\begin{theorem}[Buskin]
\label{theorem-buskin-k3-isogeny}
Let $S_1,S_2$ be projective K3 surfaces.
If
$$
H^2(S_1,\mathbb{Q})\cong
H^2(S_2,\mathbb{Q})
$$
as rational Hodge structures, then $S_1$ and
$S_2$ are symplectically isogenous.
\end{theorem}

\noindent
This is essentially a Hodge-conjecture
statement for $S_1\times S_2$.
An isomorphism of rational Hodge structures
on $H^2$ gives a rational Hodge class in
$$
H^2(S_1,\mathbb{Q})^*
\otimes H^2(S_2,\mathbb{Q})
\subset H^4(S_1\times S_2,\mathbb{Q}).
$$
To produce a symplectic isogeny, one must
realize this Hodge class by an algebraic
cycle, i.e. by a correspondence.
That is exactly the kind of assertion the
Hodge conjecture predicts.

The Shafarevich conjecture says:
if $M_1,M_2$ are hyperkahler manifolds in
the same deformation class and their rational
Hodge structures on $H^2$ are isomorphic,
then $M_1$ and $M_2$ are symplectically
isogenous.

\medskip
Very briefly: Torelli says integral periods
recover a K3 surface.
Buskin says rational periods for projective
K3 surfaces recover a symplectic isogeny.
Shafarevich asks for the analogous rational
period statement for hyperkahler manifolds.
The earlier lemma gives only the reverse
direction: an isogeny forces an isomorphism
of rational $H^2$ Hodge structures.

\begin{definition}
\label{definition-holomorphic-lagrangian-fib}
Let $(M,\Omega)$ be a holomorphically
symplectic manifold of complex dimension
$2n$.
A holomorphic map
$$
\pi:M\longrightarrow B
$$
is a {\it holomorphic Lagrangian fibration}
if its fibers have dimension $n$ and
$\Omega$ restricts to zero on the smooth
part of every fiber.
\end{definition}

\noindent
By the holomorphic Liouville theorem, the
smooth compact fibers are complex tori.
Thus a holomorphic Lagrangian fibration is a
family of complex tori over its smooth base
locus.

\medskip
{\bf Degenerate twistor deformations.}

Now let $\pi:M\to X$ be a holomorphic
Lagrangian fibration, with $M$ irreducible
holomorphic symplectic and holomorphic
symplectic form $\Omega$.
Usually $X=\mathbb{P}^n$, but this is not
essential here.
Let $\omega_{FS}$ be a Kahler form on $X$.
For $\lambda\in\mathbb{C}$, put
$$
\Omega_\lambda=
\Omega+\lambda\pi^*\omega_{FS}.
$$

\begin{lemma}
\label{claim-degenerate-twistor-form}
The closed complex form $\Omega_\lambda$
defines a complex structure on the smooth
manifold $M$ for which $\Omega_\lambda$ is a
holomorphic symplectic form.
\end{lemma}

\noindent
This is the basic degenerate twistor
deformation.
We deform the complex structure by adding a
form pulled back from the base of the
Lagrangian fibration.

\begin{definition}
\label{definition-degenerate-twistor-def}
The family of holomorphic symplectic
manifolds obtained from the forms
$$
\Omega_\lambda=
\Omega+\lambda\pi^*\omega_{FS},
\qquad \lambda\in\mathbb{C},
$$
is called the {\it degenerate twistor
deformation} associated to the Lagrangian
fibration $\pi:M\to X$.
\end{definition}

\noindent
Basic properties:
\begin{enumerate}
\item If the added form on $X$ is exact, the
deformation is trivial by Moser's theorem.
\item When $M$ is compact and of Kahler type,
these deformations are limits of ordinary
twistor deformations.
\item Each member still carries a
holomorphic Lagrangian fibration over the
same base $X$.
\item The fibers and the base are unchanged;
only the complex structure on the total
space is varied.
\item Algebro-geometrically, this is a
Tate-Shafarevich deformation: it is described
by a class in
$$
H^1(X,\operatorname{Aut}_\pi(M)),
$$
where $\operatorname{Aut}_\pi(M)$ is the
sheaf of fiberwise automorphisms.
\end{enumerate}

\begin{definition}
\label{definition-torsor}
Let $G\to X$ be a group object over $X$.
A {\it torsor} over $G$ is a space
$P\to X$ with a fiberwise action of $G$ such
that, locally on $X$, $P$ is isomorphic to
$G$ with its translation action.
Equivalently, $G$ acts freely and
transitively on the fibers of $P\to X$.
\end{definition}

\noindent
In the Lagrangian fibration setting, the
smooth fibers are tori.
Different fibrations in the same degenerate
twistor family can be viewed as torsors over
the same relative torus fibration.

\medskip
Upshot.
The culminating goal is to show that any two
projective hyperkahler manifolds in the same
degenerate twistor family are symplectically
isogenous.
To construct this isogeny, first make the two
torus fibrations literally the same after a
finite cover of the base.
Then the graph of the resulting isomorphism,
after compactification, gives the desired
correspondence.

\begin{theorem}
\label{theorem-degenerate-twistor-cover}
Let $M_1,M_2\to X$ be projective
hyperkahler manifolds in the same degenerate
twistor family.
Then there is an open subset $U\subset X$
and a finite cover $\widetilde{U}\to U$ such
that
$$
M_1\times_X\widetilde{U}
\cong
M_2\times_X\widetilde{U}
$$
over $\widetilde{U}$.
\end{theorem}

\noindent
This is the theorem where torsors enter.
After finite base change, the torsors acquire
sections, hence become trivial.
The graph of the resulting isomorphism over
an open set is then compactified to produce
the correspondence, giving the symplectic
isogeny.

\section{Hypercomplex manifolds}
\label{section-hypercomplex-manifolds}

\begin{exercise}
\label{exercise-SU2-invariant-forms}
Let $M$ be a compact hypercomplex manifold of
real dimension 4, equipped with a
quaternionic Hermitian structure, and $V$ the
space of closed
$\text{SU}(2)$-invariant 2-forms. Prove that
$V$ is finite-dimensional.
\end{exercise}

\begin{proof}[Proof by Arpan.]
Consider the Hodge star operator with respect
to the Riemannian metric on $M$.
\begin{enumerate}
\item A closed anti-self-dual form is
harmonic.
\item The space of harmonic forms is
finite-dimensional.
\item An $\text{SU}(2)$-invariant closed form
is
anti self-dual.
\begin{enumerate}
\item
$\Lambda^+=\text{span}(\omega_I,\omega_J,
\omega_K)$.
$M$ is a manifold of real dimension 4,
which means that $\dim\Lambda^2(M)=6$ and
thus
$\dim\Lambda^+(M)=3$
(here we can check using a basis of
$\mathbb{R}^6$
that $*$ permutes the basic covectors with no
fixed points,
so that it has zero trace;
in turn this means that $\dim \Lambda^+ +
\dim\Lambda^-=0$).
Further, each of $\omega_I$, $\omega_J$ and
$\omega_K$
squared and multiplied by $\frac{1}{2}$ give
the Riemannian
volume form, and therefore the three are in
$\Lambda^+(M)$.
Clearly they are linearly independent and
thus span $\Lambda^+(M)$.

\item $\omega_I,\omega_J$ and $\omega_K$ are
not fixed by the
$\text{SU}(2)$ action
Indeed, $J\omega_I=-\omega_I$ and similarly
for other cases. Therefore there are no
$\text{SU}(2)$-fixed
points in $\Lambda^+$.

\item
$\text{SU}(2)\Lambda^+(M)=\Lambda^+(M)$.
This is clear since $J\omega_I = -\omega_I$
and so on.
Since $*$ is an isomorphism of $\Lambda^+(M)$
it follows that
$\text{SU}(2)\Lambda^-(M)=\Lambda^-(M)$.

\medskip\noindent
We conclude that if $\alpha$ is
$\text{SU}(2)$-invariant (i.e. a fixed point
of
the action of $\text{SU}(2)$ on $\Lambda^2$)
its positive part will vanish, so that
$\alpha$ is anti-self-dual. More precisely,
if
$\alpha = \alpha^+ +\alpha^-$ is
$\text{SU}(2)$-invariant and $A \in
\text{SU}(2)$,
$$
\alpha^+ +\alpha^- = \alpha =
A\alpha=A\alpha^+ + A\alpha^-,
$$
so $A\alpha^+=\alpha^+$, but $\text{SU}(2)$
has no fixed points in $\Lambda^+$, so
$\alpha^+=0$.
Therefore $\alpha \in \Lambda^-$ as required.
\end{enumerate}
\end{enumerate}
\end{proof}

\begin{exercise}
\label{exercise-close-funda-forms-almos}
Let $(M,K,J,K,g)$ be an almost hypercomplex
Hermitian structure
on a K3 surface, and $\omega_I, \omega_J,
\omega_K$
its fundamental forms.

Suppse that $\omega_I$ is closed.
Prove that $\omega_J,\omega_K$ are closed, or
find a counterexample.
\end{exercise}

\noindent
The idea of the following solution
(provided entirely by ChatGPT)
is that we can ``rotate''
(in the $J,K$-plane within $\text{End}(TM)$)
two of the complex structures while
preserving the
quaternionic relations.
Any orthonormal triple of such endomorphisms
should be expected to satisfy the
quaternionic relations (algebra),
but the closedness condition on the
associated fundamental
forms will be modified (differential).

\begin{proof}
Let $(M,I,J,K,g)$ be a hyperkähler structure
on a K3 surface.
Define
\begin{align*}
\tilde{J}&=\cos \theta J+ \sin \theta K\\
\tilde{K}&=-\sin \theta J + \cos \theta K
\end{align*}

\noindent
Then for any $v \in TM$,
\begin{align*}
\tilde{J}^2v&=\cos J(\cos \theta Jv+ \sin
\theta Kv)
+\sin \theta K (\cos J v + \sin Kv)\\
&=\cos^2J^2v+\cos \theta \sin \theta J K v
+ \sin \theta \cos \theta K J v + \sin^2
\theta K^2 v\\
&=-\text{Id} v
\end{align*}

\noindent
And likewise $\tilde{K}^2=-\text{Id}$.
Also,
\begin{align*}
I \tilde{J} \tilde{K}v&=
I(\cos\theta J(-\sin \theta Jv+\cos \theta
Kv)
+\sin \theta K (- \sin Jv+ \cos \theta Kv))\\
&=I (-\cos \theta \sin \theta J^2 v
+ \cos^2 \theta J K v
- \sin^2 \theta K J v
+\sin \theta \cos \theta K^2 v)\\
&=-\cos \theta \sin \theta I J^2 v
+ \cos^2 \theta I J K
- \sin^2 \theta I K J v
+ \sin \theta \cos \theta I K^2 v\\
&=-\cos^2v- \sin^2v=-\text{Id}v.
\end{align*}

\noindent
Thus $(I,\tilde{J},\tilde{K})$ is an almost
hypercomplex
structure on $M$.
However,
\begin{align*}
d
\omega_{\tilde{J}}&=d(g(\cdot,\tilde{J}\cdot)
)\\
&=d(g(\cdot,\cos \theta J\cdot + \sin \theta
K\cdot))\\
&=d(\cos \theta \omega_J + \sin \theta
\omega_K)\\
&=d \cos \theta \wedge \omega_J + d \sin
\theta \wedge \omega_K\\
&=-\sin \theta
d\theta\wedge\omega_J+\cos\theta
d\theta\wedge\omega_K\\
&=d \theta \wedge \omega_{\tilde{K}},
\end{align*}

\noindent
which is not always zero if $\theta$ is not
constant.
\end{proof}

\section{Holonomy}
\label{section-holonomy}

\noindent
Let $M$ be a smooth manifold
and $\nabla$ a linear connection.
Any loop $\gamma$ gives a map
$P_\gamma:T_pM \to T_pM$.

Define
$$
\text{Hol}_p(\nabla)
=\{P_\gamma:\text{loops starting at }p\}
\subseteq \text{GL}(n,\mathbb{R}),
$$
which indeed is a subgroup of
$\text{GL}(n,\mathbb{R})$
since parallel transport is invertible
by changing the direction of the
parametrization.
If $M$ is connected, $\text{Hol}_p(\nabla)$
is independent of $p$.

\begin{example}
\label{example-holonomy}
\begin{enumerate}
\item For $\nabla^{\text{LC}}$
we have $\text{Hol}(\nabla^{\text{LC}})
\subset O(n,\mathbb{R})$.

\end{enumerate}
\end{example}

\begin{theorem}[Berger, 1955]
\label{theorem-berger}
If $(M,g)$ a complete Riemannian
simply connected
manifold,
not locally a product
and not a symmetric space,
then $\text{Hol}(\nabla^{\text{LC}})$
can only be one of the following:
\begin{enumerate}
\item $\text{SO}(n)$ oriented Riemannian
manifold.
\item $\text{SU}(n/2)$ Calabi-Yau.
\item $Sp(n/4)$ Hyperkähler.
\item $Sp(n/4)Sp(1)$ quaternionic Kähler.
\item $U(n/2)$ Kähler.
\item $G_2$ exceptional geometry -+.
\item $\text{Spin}(7)$ exceptional
geometry 8.
\end{enumerate}
\end{theorem}

\section{Bismut connection}
\label{section-bistmuth}

\noindent
These notes document Beatrice Brienza's
work based on a few talks by her.

\begin{theorem}[Gauduchon]
\label{theorem-gauduchon-connections}
On $TM$ there exists an affine line
of Hermitian connections known as
{\it canonical} or
{\it Gauduchon connections}
$$
g(\nabla_X^tY,Z)
=g(\nabla_X^{\text{LC}}Y,Z)
+\frac{t-1}{4}
d^c\omega(X,Y,Z)
+\frac{t+1}{4}
d^c(X,JY,JZ).
$$
\end{theorem}

\noindent
Note that:
\begin{enumerate}
\item $\nabla^{\text{LC}}$ is Hermitian
if and only if $d \omega=0$.
\item $d \omega \neq 0 \implies
\forall t \in \mathbb{R}$,
$\nabla^t$ has torsion. This leads to
``Hermitian geometry with torsion''.
\item $t=1$ gives the Chern connection,
i.e. the unique Hermitian connection
such that $\nabla^{0,1}=\bar\partial$
(see Definition
\ref{definition-chern-connection}).
\item $t=-1$ gives the Bismut connection,
defined next, satisfying
$$
g(\nabla_X^BY,Z)
=g(\nabla^{\text{LC}}_XY,Z)
-\frac{1}{2}d^c\omega(X,Y,Z),
$$
where $\omega$ is the {\it Bismut torsion}.
\end{enumerate}

\noindent
Upshot: Chern connection is adapted
for the complex structure,
Bistmut connection is adapted
for skew torsion,
both equal Levi-Civita connection in
Kähler case.

\begin{definition}[Stve, Bis, Fl, 2002]
\label{definition-bismut-connection}
The {\it Bismut connection} is the
unique Hermitian connection with
totally skew symmetric torsion, i.e.
$g(T^B(X,Y),Z) \in \Omega^3(M)$,
where of course
$T^B(X,Y)=\nabla^B_XY-\nabla^B_YX-[X,Y]$
is the usual torsion.
\end{definition}

After lowering an index of the torsion,
we obtain a tensor $T$:
$$
T \in \Lambda^2(T^*M) \otimes T^*M
\simeq \underbrace{T^*M}_{\text{vectorial
type}} \oplus\underbrace{\Lambda^3T^*M}_{
\text{skew type}}
\oplus \underbrace{\tau}_{\text{
terra incognita}}
$$
and we are interested in the case
where $T$ is in $\Lambda^3T^*M$.

\medskip\noindent
It turns out that
$$
\rho^B=0 \iff
\text{Hol}(\nabla^B) \subseteq
\text{SU}(n)
$$
where $\rho^B$ is the
Bismut Ricci curvature,
$$
\rho^B(X,Y)
=\frac{1}{2}\sum_i^{2n}
R^B(X,Y,e_i,J e_i)
\in \Omega^2(M),
$$
in fact $\rho^B \in H^2$ is
the Chern class of $M$.

\begin{definition}
\label{definition-bismut-hermite-einstein}
A {\it Bismut Hermite Einstein} manifold
is when
$$
dH=d d^c\omega=0,\qquad
\rho^B=0.
$$
\end{definition}

\noindent
And $\rho^B=0 \iff \text{Hol}(\nabla^B)
\subseteq \text{SU}(n)$.

\medskip\noindent
In dimension 4, a BHE
can only be a Kähler Ricci flat manifold
or $S^1 \times \text{SU}(2)$. On
dimension 6, ABSL.

\begin{definition}
\label{definition-almost-contact}
An {\it almost contact manifold}
is a triple $(\varphi,\eta,\xi)$
where $\varphi \in \text{End}(TM)$,
$\eta \in \Omega^1(M)$,
$\xi \in \mathfrak{X}(M)$,
$\varphi^2 = -\text{Id} + \eta \otimes \xi$,
$\eta(\xi)=1$.
\end{definition}

\noindent
This means that
$$
TM = \mathcal{F}_{\xi}\oplus Ker\xi.
$$

We then introduce a metric
compatible with $\varphi$.




\section{Instantons}
\label{section-instantons}

\noindent
Let $X$ be a compact 4-dimensional compact
oriented Riemannian manifold, $G$ a compact
Lie group (e.g. $U(n), \text{SO}(n),Sp(n)$).

\begin{definition}
\label{definition-g-instanton}
A {\it $G$-instanton} on $X$
is a $G$-vector bundle $E$
with a connection $\nabla$ such that its
curvature $F_{\nabla}$ is self-dual, i.e.
$*F_{\nabla}=F_{\nabla}$
where $*$ is Hodge operator.
\end{definition}

ADHM studied the case when $X=S^4$,
and $X=\mathbb{C}P^2$ was also studied in a
paper by N. Buchdahl.

\section{Twistor space}
\label{section-twistor-space}

\noindent

Let $X$ be a compact oriented 4-dimensional
Riemannian manifold. Recall that since the
Hodge star operator $*:\Lambda^2\to\Lambda^2$
has eigenvalues $\pm 1$, we have a
decomposition
$$
\Lambda^2=\Lambda^+ \oplus \Lambda^- \cong
\text{SO}(T_X)
$$
Consider the spherization of $\Lambda^+$:
$$
Z:=S\Lambda^+=\{S \in \Lambda^+:\text{unit
vectors}\}
$$
$(m,s) \in Z$, $m \in X$, $s_m \in
\Lambda^+_m$,

$$
T_{(m,s)}Z=
T_mX \oplus_{LC}T_sS \Lambda^+_m
\cong T\mathbb{C}P^1
$$
which has a complex structure
$\mathcal{I}_{\mathbb{C}P^1}$
(metric orientation).

\begin{theorem}
\label{theorem-twistor-space}
$(Z,\mathcal{I}_Z=(s,\mathcal{I}_{\mathbb{C}
P^1}))$
is a complex manifold of complex dimension 3.
\end{theorem}

\noindent
We also have a projection $p:Z \to X $
whose fibers are $\mathbb{C}P^1$,
called {\it real lines}.

\begin{example}
\label{example-twistor-space}
\begin{itemize}
\item $X=S^4$ then $Z=\mathbb{C}P^3$
(projective).
\item $X=\mathbb{C}P^2$ then
$Z=\mathbb{F}_3(\mathbb{C})$ (projective).
\end{itemize}
\end{example}

Let $Z$ be a 3-dimension complex manifold.
Let $L \xymatrix{\ar@{^{(}->}[r]&} Z$
be a rational curve with a normal bundle
$N_L \cong \mathcal{O}(1) \oplus
\mathcal{O}(1)$.

\begin{remark}
\label{remark-small-deformations}
If $Z \supset W$ where $W$ is a complex
submanifold,
the small deformations of $W$ are encoded by
$H^0(N_W)$ and $H^1(N_W)$.
\end{remark}

Let $\mathcal{M}$ be the space of
deformations of $L$.
Since $H^1(N_L)=0$ then
$TM \simeq H^0(N_L) = H^0(\mathcal{O}(1)
\oplus \mathcal{O}(1))$,
so $\dim_\mathbb{C}TM=4$.

Let $\sigma:Z \to Z$ be an antiholomorphic
involution with no fixed points but $L$ ($L$
is fixed). There is an induced
antiholomorphic involution $\tilde{\sigma}$
on $\mathcal{M}$.

\begin{remark}
\label{remark-antiholomorphic-involution}
\begin{itemize}
\item We can recover $X$ as a connected
component
of the fixed point set of $\tilde{\sigma}$ in
$\mathcal{M}$.
\item We can recover the
conformal structure.
\end{itemize}
\end{remark}

\medskip\noindent
Now we discuss Ward-Penrose transform.
Let $E$ be a $G$-instanton
(i.e. a vector bundle with a connection such
that
$*F_\nabla=F_\nabla$) on $X$.
Let $\phi:V \to \overline{V}^*$ be an
Hermitian form.

There is a correspondence (the {\it Ward
correspondence})
which allows to construct a holomorphic
vector bundle on $Z$
such that $\tilde{E}$ restricted to the fiber
of $p:Z \to X$
(the fibers of $p$ are $\mathbb{C}P^1$)
trivially.

Conversely, starting from a hermitian
holomorphic vector bundle $\tilde{E}$
on $Z$ which restricts trivially to the
fibers of $p:Z \to X$.

\medskip\noindent
{\bf Problem (ADHM).} Describe all
$G$-instanton bundles on $S^4$
with $r \geq 2$, $c_2=k$ (fixed second Chern
class).

This makes us want to study the vector
bundles on $\mathbb{C}P^3=Z$.

See Horrseks for the original construction of
monads.
Upshot: ``All bundles can be represented by
the cohomology of some monad''.

\begin{definition}[Monad]
\label{definition-monad}
Let $U,V,W$ be complex vector spaces over
$\mathbb{C}P^1$.
Consider the sequence $h \circ a=0$, where
$a,b$are bundle homomorphisms:
$$
\xymatrix{
0\ar[r]&U(-1):=U \otimes \mathcal{O}(1)
\ar[r]^a
& V \otimes \mathcal{O}
\ar[r]^b
& W(1):=W \otimes \mathcal{O}(1)\ar[r]&0
}
$$
Let $E:=\frac{\Ker b}{\text{Im}a}$.
\end{definition}

\medskip\noindent
Let's study the case of $\mathbb{C}P^3$.
If $Z_i$ are homogeneous coordinates on
$\mathbb{C}P^3$,
$$
a=\sum_i A_iZ_i,\qquad b=\sum_jB_jZ_j
$$
the {\it monad condition} is
$$
B_jA_i+A_jB_j=0
$$
[which gives the conditions for a hyperkähler
metric on the moduli space].

\begin{theorem}[Barth, Barth-Hulek]
\label{theorem-bbh}
Let $E$ be a holomorphic vector bundle of
rank $r \geq 2$
on $\mathbb{C}P^3$ with a $c_2=k$ and
$E|_L=$trivially.
Then $E$ is the cohomology of a monad.
Also, it satisfies some vanishing conditions:
$$
H^0(E)=H^0(E^*)=H^1(E(-2))=H^1(E^*(-2))=0.
$$
\end{theorem}






\section{Stability}
\label{section-stability}

\noindent
This section outlines some concepts
from a paper by Verbitsky from approx. 2002.

A {\it Gauduchon metric} is such that
$\partial\bar\partial \omega^{n-1}=0$. The
degree of a holomorphic vector bundle is
given by $$ \text{deg}_g(\mathcal{F}) =
\int_M c_1(\det \mathcal{F}) \wedge
\omega^{n-1}, $$ which is well defined when
$g$ is Gauduchon.

\begin{definition}
\label{definition-slope-and-semistability}
The {\it slope} of a vector bundle is
$$
\mu_g(\mathcal{F}) = \frac{\text{deg}_g
\mathcal{F}}{\text{rk}\mathcal{F}}.
$$
A vector bundle $\mathcal{F}$ is {\it stable}
if for any coherent subsheaf $\mathcal{E}
\subset \mathbb{F}$
$$
\mu_g(\mathcal{E}) < \mu_g(\mathcal{F}).
$$
It is called {\it semistable} if we put $\le
$
instead.
\end{definition}

\begin{example}
\label{example-line-bundles-are-stable}
Line bundles are stable.
\end{example}

\begin{theorem}[Hitchin-Kobayashi
correspondence]
\label{theorem-hitchin-kobayashi}
Let $(M,g)$ be a compact complex manifold
with $g$ Gauduchon.
Then an hermitian vector bundle $(F,h)$
is Hermitian-Einstein,
i.e. it satisfies
$$
\sqrt{-1}\Lambda\Theta=\gamma\text{id}
$$
if and only if
it is polystable, i.e.
it is the sum of stable bundles
with the same slope.
\end{theorem}

\noindent
The aim of Misha's paper is to
classify stable bundles in positive elliptic
fibrations using Hitchin-Kobayashi
correspondence. Positive elliptic fibrations
are $\pi:M \to X$ with an action of a torus
on the fibers, and such that
$\pi^*\omega_X=d\eta$ where $\omega_X$ is a
Kähler form on $X$. Such a metric is called
preferred if it is invariant on the fibres
and is a Riemannian submersion.

The main theorem is the following:

\begin{theorem}
\label{theorem-stability-misha}
For a positive elliptic fibration $\pi:M \to
X$,
if $(B,h)$ is a stable bundle on $M$,
$$
B \simeq L \otimes \pi^*B_0
$$
where $L$ is a line bundle and
$B_0$ is stable bundle on $X$.
\end{theorem}

\noindent
The upshot is that a locally conformal
hyperkähler manifold
has an associated elliptic fibration…


\section{Deformation theory (introduction)}
\label{section-deformations}

\noindent
Here I discuss deformations in the category
of complex spaces.

Let's just do a summary of \cite[Chapter
9]{Voisin}.

Let $\mathcal{X}$ be a complex manifold, $B$
a complex manifold
and $\phi:\mathcal{X} \to B$ a holomorphic
map.
We call $\phi$ a {\it family} if it is a
proper holomorphic submersion.

A theorem by Ehresmann says that any
``deformation'' of
smooth real manifolds is trivial:

\begin{theorem}[Ehresmann]
\label{theorem-ehresmann}
\begin{reference}
\cite[Theorem 9.3]{voi}
\end{reference}
If $\phi:\mathcal{X} \to B$ is a
proper submersion of differentiable manifolds
(no complex structure), with $B$ contractible
and $0 \in B$
some distinguished point,
then the total space $\mathcal{X}$ is trivial
as a
manifold over $B$, i.e.
there exists a diffeomorphism
 $$
T:\mathcal{X}\xrightarrow{\cong}X_0 \times B
$$
of manifolds over $B$, i.e. commuting with
projection.
\end{theorem}

\begin{proof}
The main tool is the theorem of tubular
neighbourhoods,
which says that there is a diffeomorphism of
a nieghbourhood $W$ of $X_0$ in $\mathcal{X}$
and a neighbourhood of $X_0$ in its normal
bundle.
(We have a normal bundle because $X_0$ is a
submanifold of $\mathcal{X}$.)

In particular this gives a differentiable
retraction
$T_0:W \to X_0$, which allows us to consider
the map
$$
(T_0,\phi):W \to X_0 \times B.
$$
Since $T_0$ is a retraction
(it's the identity when restricted to $X_0$)
and $\phi$ is a submersion,
the map $(T_0,\phi)$ has nonsingular
differential along $X_0$.
Then we can invert it locally along $X_0$,
and since $X_0$ is compact we obtain a whole
neighbourhood $W'\subset W \subset
\mathcal{X}$
where the differential is nonzero and thus
$(T_0,\phi)|_{W'}$ is an embedding
(using the inverse function theorem).

Finally we restrict the image of
$(T_0,\phi)|_{W'}$
to a neighbourhood $U\subset B$ of $0$,
using properness of $\phi$ to ensure
that $\phi^{-1}(U) \subset W'$ --- though I'm
not exactly sure how.
We have constructed a diffeomorphism
$\phi^{-1}(U) \to X_0 \times U$
commuting with the projection.

This proof was only local; the global version
is in \cite{voi} though it's not used in
the rest of the text.
\end{proof}

In the complex situation, if
$\phi:\mathcal{X} \to B$ is a family
of complex manifolds and $0 \in B$, then
``up to replacing $B$ by a neighbourhood of
0'',
there exists a smooth trivialization
$T=(T_0,\phi):\mathcal{X} \to X_0 \times B$
such that the fibres of $T_0$ are complex
submanifolds of $\mathcal{X}$.

Now let $\phi:\mathcal{X} \to B$ be a family
of complex manifolds.
Recall the pullback bundle $\phi^*B$ and the
fact
that the differential $\phi_*$ can be thought
of as a gadget
that maps tangent vectors to $\mathcal{X}$
not to vectors tangent to $B$ but
to vectors in the pullback bundle $\phi^*B$.
Restricting everything to the fibre
$X=\phi^{-1}(0)$ we get
\begin{equation}
\label{equation-short-exact-seque-defor}
\xymatrix{
0\ar[r]&T_X\ar[r]&T_{\mathcal{X}|X}\ar[r]&
\phi^*T_{B|X}\ar[r]&0.
}
\end{equation}
There's what seems to be a really nice result
saying that
$H^1(\mathcal{F})$ parametrizes isomorphisms
classes
of extensions of $\mathcal{F}$ by the trivial
bundle,
but for now let me just define it roughly as
follows.
The category of $\mathcal{O}_X$ is most
surely abelian
with enough injectives. This means that
objects
have associated complexes, and we can take
cohomologies of these complexes.
This says, roughly, that short exact
sequences of vector bundles have associated
cohomology
long exact sequences. Then the
Kodaira-Spencer map
is the map
$$
\rho:T_{B,0}=H^0(X,\phi^* T_{B|X})\to
H^1(X,T_X)
$$
induced by the long exact sequence associated
to
\ref{equation-short-exact-seque-defor}.
Here the equality $T_{B,0}=H^1(X,\phi^*
T_{B|X})$
is just because the pullback of the tangent
bundle of $B$
to any fiber is trivial.

\medskip\noindent
Now I will explain first order deformations.
Consider a family $\mathcal{X} \to B$.
We introduce a complex analytic subscheme
$B_\varepsilon$ of $B$
giving an ideal sheaf $\mathcal{M}^2_0$
where $\mathcal{M}_0$ is the maximal ideal
of $\mathcal{O}_{B,0}$ consisting of the
holomorphic
functions vanishing at 0.
This means that a function in
$\mathcal{M}^2_0$
is (a finite sum of terms) of the form $f=gh$
for $g,h \in \mathcal{M}_0$,
which means that $df=d(gh)=hdg+gdh=0$ at $0$
since both $g$ and $h$ vanish at $0$.
That is, this sheaf contains functions that
vanish
at $0$ and whose first derivative vanishes at
$0$.

This means that $B_\varepsilon$ is a point:
it's the vanishing locus of the set of
functions
that vanish at a single point.
However, it is endowed with an unusual
structure sheaf.
Let $v$ be a tangent vector to $B$ at $0$ and
$f$ a function in the strange sheaf
$\mathcal{M}_0/\mathcal{M}_0^2$.
Then by the computation above, since $v$ is a
differential
operator, $v$ annihilates $f$.
This means that this sheaf detects not only
the point $0$
but also any tangent vector at $0$.
Let's denote
$B_\varepsilon=\Spec
\mathbb{C}[\varepsilon]/(\varepsilon^2)$
just because its structure sheaf is
$\mathbb{C}[\varepsilon]/(\varepsilon^2)$.

In fact, this is so throughout deformation
theory:
$\Spec k[\varepsilon]/(\varepsilon^2)$ is
topologically a point!
But it's a {\it fat point}: the support of
its
structure sheaf is a point, but as a scheme,
it's not a point.
And the same happens with the family
$\mathcal{X}$ over
$\Spec k[\varepsilon]/(\varepsilon^2)$,
topologically it's just the fibre $X$!
But its structure sheaf is different:
``it contains the directions in $\mathcal{X}$
which
are normal to $X$''.

\medskip\noindent
Perhaps the most important part of this
construction
comes when considering a fat point
$B_\varepsilon \subset B$.
This gives a first-order deformation
$\phi^{-1}(B_\varepsilon)) \subset
\mathcal{X}$
$$
\xymatrix{
X\ar[r]\ar[d]
&\phi^{-1}(B_\varepsilon) \subset
\mathcal{X}\ar[d]\\
\ast\ar[r]
&B_\varepsilon \subset B
}
$$
Using the constructions above, one shows that
to any such deformation we can associate a
cocycle
in $H^1(X,T_X)$.
Thus, first order deformations are ``exactly
parametrised''
by the group $H^1(X,T_X)$.
This is the main idea behind Kodaira-Spencer
map,
which is well-explained in
\href{https://en.wikipedia.org/wiki/Kodaira
% E2%80%93Spencer_map}{Wikipedia}.

I will present a similar construction
to Kodaira-Spencer map
that works for embedded deformations
instead of just deformations.
(That's with the intention of understanding
the
forthcoming exercise Exercise
\ref{exercise-deformations-curve-in-k3}.)

Since we are in complex space category
I follow \cite[p. 224]{voi} for the main
setup,
but this is also a mixture
of \cite[p. 147, 150]{Sernesi}.

The result is in fact quite reasonable
geometrically:
the dimension of directions in which we can
deform a
submanifold should be non other than
the dimension of the normal bundle.

\section{Deformation theory}
\label{section-deformation-theory}

\begin{definition}
\label{definition-deformation}
Let $X$ be a complex manifold.
A {\it deformation} of $X$ is a holomorphic
submersion
$\pi:\mathcal{X} \to B$ of complex manifolds
such that $\pi^{-1}(0)\cong X$ for some $0
\in B$.

If $\mathcal{X}$ and $\mathcal{Y}$ are two
deformations
of $X$ over the same base,
a {\it morphism of deformations} is a map
$f:\mathcal{X} \to \mathcal{Y}$ such that the
diagram
$$
\xymatrix{
& X\ar[dl]\ar[dr]\\
\mathcal{X}\ar[dr]\ar[rr]^f&&\mathcal{Y}
\ar[dl]\\
&B
}
$$
commutes. The {\it space of deformations} of
$X$
is the set of deformations modulo
the equivalence relation of two deformations
being isomorphic. A priori we don't know
if this set has any structure.

Let $Y$ be another complex manifold.
A {\it deformation of $X$ in $Y$}
is a deformation of $X$ such that
$\mathcal{X}\subset Y\times B$.

The {\it fat point} is a ringed space
$(B_\varepsilon,\mathbb{C}[\varepsilon]/
(\varepsilon^2))$
whose underlying topological space
$B_\varepsilon$ is a point
and structure sheaf is the ring of dual
numbers
$\mathbb{C}[\varepsilon]/(\varepsilon^2)$.

A {\it first order deformation} of $X$ is a
deformation such that
the base $B$ is a fat point.
\end{definition}

The following proposition is a
mixture of ideas from \cite[p. 224]{voi}
and \cite[Examples
3.2.4(ii)]{Sernesi-deformations}.

\begin{proposition}
\label{proposition-kodai-spenc-embed-defor}
Let $X$ be a smooth hypersurface in a complex
manifold $Y$.
For every deformation of $X$ in $Y$
with base of dimension 1 we can construct a
section of the bundle
$N_{X/Y}$. Different deformations give
different sections up to rescaling.
\end{proposition}

\begin{proof}
Let
$$
\xymatrix{
X\ar[r]\ar[d]
&\mathcal{X}\subset Y\times B\ar[d]^\pi\\
0\ar[r]
&B
}
$$
be a deformation of $X$ in $Y$
and suppose that $\dim B=1$.
This means that $\mathcal{X}$ is a complex
submanifold of
the product $Y \times B$
and $\pi:\mathcal{X} \to B_\varepsilon$ is
a holomorphic submersion.

Since $X$ is a hypersurface,
we can find an open cover $V_i$
of $M$ such that $V_i \cap X:= U_i$
is given locally by the vanishing of
a function $f_i \in \mathcal{O}_Y(V_i)$.
In terms of coordinates we find that
that $V_i \cap X$ is given by
$(z_1,…,z_{n-1},0)$.

Since $\pi$ is a submersion,
by the holomorphic inverse function theorem
we obtain for every $V_i \subset \mathcal{X}$
(such that $V_i \cap X \neq  \emptyset$)
a biholomorphism from some $B_i \subset B$
containing $0$ to $V_i$.
This gives coordinates $(0,…,b)$ to $V_i$.
We conclude that $V_i \cong U_i\times B_i$
along $X$.

Now we ``restrict'' to a first order
deformation.
Consider the fat point subscheme
$B_\varepsilon$ of $B$
(I think this is the reason why $B$ must have
dimension 1,
in Voisin notation, that a first order
neighbourhood
of $0$ in $B$ is indeed a fat point)
and let
$X_\varepsilon=\pi^{-1}(B_\varepsilon)$,
which is a ringed space whose topological
space
is the same as $X$, but its
structure sheaf is induced from
$B_\varepsilon$.

(I think that to prove this part I have
to look carefully at the induced sheaf
$\pi^*(\mathbb{C}[\varepsilon]/
(\varepsilon^2)$.)
Since we are on a local trivialization,
the structure sheaf of the product $U_i
\times B_i$
is the tensor product of the structure
sheaves
(reference?),
so that the structure sheaf of
$X_\varepsilon$ at
$U_i\times B_\varepsilon$ is given by
$$
\mathcal{O}_{X_\varepsilon}(U_i)=
 \mathcal{O}_Y(U_i) \otimes
 \mathbb{C}[\varepsilon]/(\varepsilon^2).
$$
Functions here look like finite sums of
functions of the form $f+ \varepsilon g$ for
$f,g \in \Gamma(U_i,\mathcal{O}_{U_i})$.

The next key step is to argue that
$X_\varepsilon$ is locally cut by
functions of the kind $f_i + \varepsilon g_i$
for some $g_i \in \mathcal{O}_M(U_i)$.
Here the $f_i$ are the same as those which
define $X$.
I think a proof is that from the ideal sheaf
exact sequence of $X$ in $Y$
$$
\xymatrix{
0\ar[r]
&\mathcal{I}_X\ar[r]
&\mathcal{O}_Y\ar[r]
&\mathcal{O}_X\ar[r]
&0
}
$$
we obtain after tensoring by
$\mathbb{C}[\varepsilon]/(\varepsilon^2)$
$$
\xymatrix{
0\ar[r]
&\mathcal{I}_X \otimes
\mathbb{C}[\varepsilon]/(\varepsilon^2)\ar[r]
&\mathcal{O}_Y \otimes
\mathbb{C}[\varepsilon]/(\varepsilon^2)\ar[r]
&\mathcal{O}_X \otimes
\mathbb{C}[\varepsilon]/(\varepsilon^2)\ar[r]
&0.
}
$$
(But is the functor $-\otimes
\mathbb{C}[\varepsilon]/(\varepsilon^2)$
exact? Why? So far we didn't impose a
flatness condition on $\pi$.)
The sheaf in the right
coincides with the structure sheaf of
$X_\varepsilon$ when
evaluated on a trivial open set of the form
$U_i \otimes B_\varepsilon$.
This shows that the ideal sheaf of
$X_\varepsilon$
as a ringed space is
$\mathcal{I}_X
\otimes\mathbb{C}[\varepsilon]/
(\varepsilon^3)$.
Since $\mathcal{I}_X$ is locally principal
generated by $f_i$,
we conclude that indeed $X_\varepsilon$
must be cut locally by functions of the form
$f_i + \varepsilon g_i$
for some $g_i \in \mathcal{O}_Y(U_i)$

Recall that the line bundle associated to $X$
is $\mathcal{O}_Y(X)$, given by the cocycle
$f_i/f_j:=f_{ij}$.
A version of adjunction formula
(more precisely, \cite[Adjunction formula I,
p. 146]{GH})
establishes that this line bundle,
when restricted to $C$, is in fact
the normal bundle of $X$ in $Y$, defined as
que quotient
bundle $T_Y'/T_X':=N_{X/Y}$.
In symbols, $\mathcal{O}_Y(X)|_{X}\cong
N_{X/Y}$.

Similarly, the line bundle associated to
$\mathcal{X}$
is given by the cocycle
$$
F_{ij}=f_{ij}+\varepsilon g_{ij}
$$
for $g_{ij}=\frac{g_i-f_{ij}g_j}{f_j}$.
Indeed, note that
$\frac{1}{f_j+\varepsilon
g_j}=\frac{1}{f_j}-\varepsilon
\frac{g_j}{f^2_j}$
and compute $F_{ij}=(f_i+\varepsilon
g_j)/(f_j+\varepsilon g_j)$.

Then
$$
f_i+\varepsilon g_i=(f_{ij}+\varepsilon
g_{ij})(f_j+\varepsilon g_j)
$$
gives
$$
g_i=f_{ij}g_j+g_{ij}f_j
$$
and since $f_j$ vanishes along $X$,
we are left with $g_i=f_{ij}g_j$
which is how a section of the sheaf $N_{X/Y}$
is obtained.
\end{proof}

\begin{exercise}
\label{exercise-deformations-curve-in-k3}
Let $C$ be a smooth genus $g$ curve which
can be embedded in a K3 surface $M$, and
$\mathcal{X}$
``the family of all deformations'' of $C$ in
$M$.
Prove that $\dim \mathcal{X} \leq g$.
\end{exercise}

\begin{proof}
By Adjunction formula II
(Theorem \ref{theorem-adjunction-ii})
and $M$ being a K3 surface
we get that $h^0(C,N_{C,M})=g$.

If the deformation space $X$ exists and is a
complex manifold,
the map constructed in
Proposition
\ref{proposition-kodai-spenc-embed-defor}
would be a holomorphic (why holomorphic!?)
injective map $X \to H^0(C,N_{C/Y})$.
This obliges $\dim X \leq  h^0(C,N_{C/Y})=g$.
Indeed, if not, the differential must have a
kernel,
and integrating along a direction in the
kernel
we obtain the germ of a complex curve along
which
the map is constant, which is impossible
since it is injective.
\end{proof}




\bibliography{my}
\bibliographystyle{amsalpha}



\end{document}
